\documentclass[11pt]{amsart}
\usepackage[T1]{fontenc}
\usepackage{lmodern,microtype,amsmath,amssymb,amsthm,mathtools,booktabs,enumitem}
\usepackage[margin=1in]{geometry}
\usepackage[hidelinks]{hyperref}
\newtheorem{theorem}{Theorem}[section]
\newtheorem{proposition}[theorem]{Proposition}
\newtheorem{lemma}[theorem]{Lemma}
\newtheorem{corollary}[theorem]{Corollary}
\theoremstyle{definition}
\newtheorem{definition}[theorem]{Definition}
\theoremstyle{remark}
\newtheorem{remark}[theorem]{Remark}
\newcommand{\bits}{\{-1,1\}}
\newcommand{\E}{\mathbb E}
\renewcommand{\P}{\mathbb P}
\newcommand{\HH}{\mathsf H}
\newcommand{\I}{\mathsf I}
\newcommand{\D}{\mathsf D}
\newcommand{\wh}[1]{\widehat{#1}}
\newcommand{\eps}{\varepsilon}
\newcommand{\supp}{\operatorname{supp}}
\newcommand{\Var}{\operatorname{Var}}
\newcommand{\Maj}{\operatorname{Maj}}
\newcommand{\TC}{\operatorname{TC}}
\newcommand{\KL}{\operatorname{KL}}
\DeclareMathOperator{\atanh}{atanh}
\setlength{\emergencystretch}{2em}
\setlist{itemsep=3pt,topsep=5pt}
\raggedbottom
\title[Restrictions and witness coding for FEI]{Restrictions and witness coding for the Fourier entropy--influence problem}
\author{Yitzchak Shmalo}
\address{Einstein Institute of Mathematics, The Hebrew University of Jerusalem, Jerusalem, Israel}
\email{yitzchak.shmalo@gmail.com}
\date{11 October 2026}
\subjclass[2020]{42C10, 68Q17, 60C05, 94B05}
\keywords{Fourier entropy, total influence, Boolean function, restriction, spectral coding, composition}
\hypersetup{pdftitle={Restrictions and witness coding for the Fourier entropy-influence problem},pdfauthor={Yitzchak Shmalo}}
\begin{document}
\begin{abstract}
We study reductions and extensions of the Fourier entropy--influence conjecture. Its optimal constant is preserved under simultaneous restrictions to nearly constant functions, arbitrarily small total and relative coordinate influences, and arbitrarily separated Fourier supports. These examples retain their spectral entropy under prescribed finite depths of coordinate restriction. We refine a witness-forest entropy bound by conditioning on nonempty roots and subtracting posterior witness entropy, and prove a composition bound for multiplexers with shared data variables. Exact restriction identities and biased monotone examples identify limitations of local entropy charges and threshold normalization estimates. The results leave a concrete universal witness-cost inequality open and do not establish the unrestricted conjecture.
\end{abstract}
\maketitle

\section{Introduction}

For a Boolean function $f:\bits^n\to\bits$, its squared Fourier coefficients form a probability law
\[
 P_f(S)=\wh f(S)^2,\qquad
 f=\sum_{S\subseteq[n]}\wh f(S)\chi_S,
 \qquad \chi_S(x)=\prod_{i\in S}x_i.
\]
Write
\begin{equation}\label{eq:HI}
 \HH(f)=\sum_S P_f(S)\log\frac1{P_f(S)},\qquad
 \I(f)=\sum_S |S|P_f(S),\qquad
 \I_i(f)=\sum_{S\ni i}P_f(S).
\end{equation}
All logarithms denoted by $\log$ are to base $2$; $\ln$ denotes the natural logarithm. The convention $0\log(1/0)=0$ is used throughout. The Fourier entropy--influence conjecture of Friedgut and Kalai~\cite{FK96} asserts the existence of a universal finite $C$ such that $\HH(f)\le C\I(f)$.

The companion threshold manuscript~\cite{Threshold} proves this estimate for functions with an ordered desirability relation, including all linear threshold functions, and gives extensions to combinations and unions of such functions. Its bounds for overlapping gates retain a logarithmic dependence on the overlap. Its product-measure composition result covers symmetric threshold gates. These assumptions are substantial: representing every Boolean function by a formula does not by itself control the entropy cost of shared variables.

The present results address three distinct issues. First, several superficially restrictive classes still retain the full optimal FEI constant, even when imposed simultaneously. Second, a forest reference can charge a gate name once for an entire Fourier set and recover the entropy of latent choices. This improves the coordinatewise overlap estimate on addressing functions. Third, exact composition and restriction calculations separate valid closure statements from helper inequalities that fail.

The coding and activation arguments build on earlier internal notes~\cite{BrainErasure,BrainCritical}. Small-influence reductions also have a public antecedent: Shalev~\cite{Shalev18} proves FEI when $\I(f)\le2^{-cn}$ for fixed $c>0$ and shows that an extension to $\I(f)\le2^{-s(n)}$ with $s(n)=o(n)$ would imply unrestricted FEI. The combined reduction below adds a variance-weighted bound for the redundancy cost and smooths only the auxiliary coordinates, so that small relative influences and support separation can be imposed at the same time. The composition framework is related to O'Donnell and Tan~\cite{ODT13}; the multiplexer statements below are proved directly. The restriction calculations are related to the stopping-tree approach of Gonz\'alez, MacManus and Pereyra~\cite{GMP26}. None of the results below bounds the unrestricted optimal constant.

The distinction between a class theorem and a universal reduction also applies to percolation. The entropy theorem in~\cite{Percolation} proves an FEI estimate for planar crossing functions using spatial structure and arm estimates. Those estimates furnish a geometric way to pay an encoding cost. They do not supply an encoding theorem for arbitrary Boolean spectra. The general entropy estimate of Kelman, Kindler, Lifshitz, Minzer and Safra~\cite{KKLMS20} instead has an additional logarithmic degree weight. Removing such a weight requires an additional argument.

\section{Simultaneous reductions preserving the optimal constant}\label{sec:hardening}

Put
\[
 C_{\rm FEI}=\sup_{f:\,\I(f)>0}\frac{\HH(f)}{\I(f)}\in[0,\infty].
\]
For a family of subsets $\mathcal S$, its distance is the minimum of $|S\mathbin\triangle T|$ over distinct $S,T\in\mathcal S$. A singleton family has infinite distance. A coordinate restriction fixes physical input coordinates to signs. It does not condition the spectral law.

\begin{theorem}[Simultaneous reduction]\label{thm:simultaneous}
Fix $\eps>0$ and an integer $k\ge0$. Consider nonconstant Boolean functions $F$ satisfying
\begin{equation}\label{eq:hardclass}
 \I(F)\le\eps,\qquad
 \E F\ge1-\eps,\qquad
 \max_i\frac{\I_i(F)}{\Var(F)}\le\eps,
 \qquad d(\supp\wh F)>k.
\end{equation}
The supremum of $\HH(F)/\I(F)$ over this class is exactly $C_{\rm FEI}$, including when this supremum is infinite.

More strongly, for every nonconstant seed $f$ and $\delta>0$, there exists such an $F$ with $d(\supp\wh F)>2k$ and
\[
 \frac{\HH(F)}{\I(F)}\ge\frac{\HH(f)}{\I(f)}-\delta,
\]
such that every restriction fixing at most $k$ coordinates preserves its entropy and mean. Every such restriction satisfies~\eqref{eq:hardclass} and has at least the same entropy-to-influence ratio as $F$.
\end{theorem}

The theorem does not assert that these functions violate FEI. It says that proving FEI for the class in~\eqref{eq:hardclass} is already sufficient for the full conjecture, with no loss in the optimal constant. Several preliminary identities give the proof.

\subsection{Balanced composition and a diffuse seed}

\begin{lemma}[Balanced substitution]\label{lem:balanced}
Let $u$ be Boolean on $N$ coordinates and let $g_1,\ldots,g_N$ be balanced Boolean functions on disjoint variable blocks. Then $U=u(g_1,\ldots,g_N)$ satisfies
\begin{align}
 \E U&=\E u,\label{eq:compmean}\\
 \HH(U)&=\HH(u)+\sum_i\I_i(u)\HH(g_i),\label{eq:comph}\\
 \I(U)&=\sum_i\I_i(u)\I(g_i).\label{eq:compi}
\end{align}
For a leaf $j$ in block $i$, $\I_j(U)=\I_i(u)\I_j(g_i)$.
\end{lemma}
\begin{proof}
Expand $u$ in characters and each $g_i$ in its Fourier basis. Since $g_i$ is balanced, its spectral sets are nonempty. The set of occupied blocks therefore uniquely identifies the outer spectral set. Conditional on that set, the inner spectral sets are independent and have the laws $P_{g_i}$. The chain rule for entropy and the expected size of the resulting set give the formulas.
\end{proof}

\begin{lemma}[Balanced diffuse seeds]\label{lem:diffuse}
For every nonconstant Boolean $f$ and positive $\eta,\delta$, there is a balanced Boolean $b$ with
\[
 \max_i\I_i(b)\le\eta,
 \qquad \frac{\HH(b)}{\I(b)}\ge\frac{\HH(f)}{\I(f)}-\delta.
\]
\end{lemma}
\begin{proof}
Multiplying $f$ by an appropriate character preserves its entropy and changes each influence to $\min(\I_i(f),1-\I_i(f))$. Call the resulting function $u$, and set $J=\I(u)\le\I(f)$. If $\HH(f)=0$, any sufficiently deep balanced majority composition has the required properties. Otherwise $J>0$.

On disjoint variables take $b_0=\Maj_3\prod_{j=1}^q u_j$, where the $u_j$ are copies of $u$. Then
\[
 \E b_0=0,\quad L:=\I(b_0)=qJ+\tfrac32,
 \quad H_0:=\HH(b_0)=q\HH(f)+2,
 \quad\max_i\I_i(b_0)\le\tfrac12.
\]
For the balanced depth-$d$ self-composition $b_d$, Lemma~\ref{lem:balanced} gives
\[
 \I(b_d)=L^d,\qquad
 \HH(b_d)=H_0\frac{L^d-1}{L-1},\qquad
 \max_i\I_i(b_d)\le2^{-d}.
\]
First choose $q$ large and then $d$ large. The ratio approaches $\HH(f)/J\ge\HH(f)/\I(f)$, while the maximum coordinate influence approaches zero.
\end{proof}

\subsection{Sparse activation}

Let $h$ be Boolean with mean $a$, and let $A$ be the indicator of one vertex of an independent $s$-dimensional control cube. Thus $p:=\E A=2^{-s}$. Define
\begin{equation}\label{eq:activate}
 G=1+A(h-1),\qquad d=1-a.
\end{equation}
The function equals $h$ on the active control fiber and equals $1$ elsewhere.

\begin{lemma}[Activation identities]\label{lem:activation}
The function in~\eqref{eq:activate} satisfies
\begin{equation}\label{eq:activationI}
 \E G=1-pd,\qquad \Var(G)=2pd-p^2d^2\le4p,
 \qquad \I(G)=p\I(h)+sp d.
\end{equation}
Writing $\phi(t)=-t\log t$ and $\HH_{\ne\varnothing}(h)=\HH(h)-\phi(a^2)$, its entropy is exactly
\begin{align}
 \HH(G)={}&p\HH_{\ne\varnothing}(h)+2sp(1-a^2)\notag\\
 &+(p-p^2)d^2\bigl(2s-2\log d\bigr)
   +\phi\bigl((1-pd)^2\bigr),\label{eq:activationH}
\end{align}
where the term involving $d^2\log d$ is interpreted continuously at $d=0$.
\end{lemma}
\begin{proof}
Choose the active vertex to be all $+1$, so that every Fourier coefficient of $A$ is $p$. For a nonempty data set $S$ and any control set $T$, the coefficient of $G$ is $p\wh h(S)$. For $S=\varnothing$ and $T\ne\varnothing$, it is $p(a-1)$. The empty coefficient is $1-pd$. These three disjoint groups give~\eqref{eq:activationH}. A data influence scales by $p$. Each control influence is $2p\P(h=-1)=pd$. This proves~\eqref{eq:activationI}.
\end{proof}

In particular, if $b$ is balanced and $h=b^{\otimes m}$ is a product of $m$ independent copies, then
\begin{equation}\label{eq:balancedactivation}
 \I(G)=pm\I(b)+sp,\qquad
 \HH(G)=pm\HH(b)+O(ps),\qquad
 \Var(G)=2p-p^2.
\end{equation}
The entropy error has an absolute constant for $s\ge1$. Choosing $s=\lceil2\log m\rceil$ makes $\I(G)$ tend to zero while preserving the seed ratio asymptotically.

\subsection{Systematic spectral coding}

\begin{lemma}[Variance-weighted coding]\label{lem:coding}
Let $G$ be Boolean on $N$ coordinates and let $A$ be an $r\times N$ binary matrix. For $z\in\bits^r$, put
\[
 F(x,z)=G\bigl(x_i\chi_{A_i}(z):1\le i\le N\bigr),
\]
where $A_i$ is the $i$th column. Then
\begin{align}
 \E F&=\E G,\qquad \HH(F)=\HH(G),\notag\\
 \I(F)&=\I(G)+\sum_S |AS|P_G(S)
 \le\I(G)+r\Var(G).\label{eq:codecost}
\end{align}
The original coordinate influences are unchanged, and each added coordinate has influence at most $\Var(G)$.
\end{lemma}
\begin{proof}
The Fourier expansion is
\[
 F(x,z)=\sum_S\wh G(S)\chi_S(x)\chi_{AS}(z).
\]
The map $S\mapsto(S,AS)$ is injective, so it preserves all coefficient magnitudes and the empty coefficient. Its extra degree is $|AS|$. Since $A\varnothing=0$, only nonempty spectral mass contributes to this extra degree; that mass is $\Var(G)$. This proves the sharper last bound in~\eqref{eq:codecost}.
\end{proof}

\begin{lemma}[A short redundancy bound]\label{lem:codeexist}
For fixed integer $d_0\ge1$ and sufficiently large $N$, there is an $r\times N$ binary matrix with
\[
 r=3d_0\lceil\log(2N)\rceil
\]
such that the systematic code $\{(S,AS):S\subseteq[N]\}$ has distance greater than $d_0$.
\end{lemma}
\begin{proof}
Choose the entries of $A$ independently and uniformly. For each nonzero set $S$, $AS$ is uniform on $\{0,1\}^r$. A union bound for nonzero codewords of weight at most $d_0$ gives
\[
 2^{-r}\sum_{j=1}^{d_0}\binom Nj
       \sum_{\ell=0}^{d_0-j}\binom r\ell
 \le2^{-r}d_0(N+r)^{d_0}.
\]
For sufficiently large $N$, $r\le N$, and the last expression is at most $d_0/(2N)^{2d_0}<1$. Hence a suitable matrix exists.
\end{proof}

\subsection{Smoothing the auxiliary coordinates}

The coded function from Lemma~\ref{lem:coding} may have added influences comparable to its variance. This requires a further step when proving Theorem~\ref{thm:simultaneous}.

For odd $L$ and positive integer $t$, let
\[
 g=\Maj_L\bigl(\chi_{[t]}(z^{(1)}),\ldots,
                  \chi_{[t]}(z^{(L)})\bigr),
\]
on disjoint parity blocks. This balanced gadget has
\begin{equation}\label{eq:gadget}
 \I(g)=t\I(\Maj_L),\quad \HH(g)=\HH(\Maj_L),\quad
 \max_j\I_j(g)=\beta_L:=\binom{L-1}{(L-1)/2}2^{-(L-1)}\longrightarrow0.
\end{equation}
Its nonempty Fourier supports have size at least $t$. Distinct supports have distance at least $2t$: majority has only odd Fourier degrees, and two distinct odd subsets differ in at least two positions, each replaced by an entire $t$-block.

\begin{proof}[Proof of Theorem~\ref{thm:simultaneous}]
Choose, by Lemma~\ref{lem:diffuse}, a balanced seed $b$ with $\max_i\I_i(b)\le\eps$ and ratio at least that of $f$ minus $\delta/2$. Its number $n$ of variables is fixed during the following limit. Choose an odd $L$ so that $\beta_L\le\eps$, and put $t=2k+1$ and $d_0=2k+1$. These choices are fixed as well.

Activate $h=b^{\otimes m}$ with $s=\lceil2\log m\rceil$ controls and $p=2^{-s}$. Code the resulting $N=mn+s$ coordinates with the systematic code of Lemma~\ref{lem:codeexist}, using $r=O_{k,n}(\log m)$ additional coordinates. Call the coded function $F_0$.

Now replace only the $s$ activation controls and the $r$ new parity coordinates of $F_0$ by disjoint copies of the gadget $g$. Leave every data coordinate unchanged. By balanced substitution, the mean and variance are unchanged. The total influence in the coordinates being replaced is at most
\[
 sp+r\Var(G)\le p(s+2r),
\]
because $b$ is balanced and $\Var(G)=2p-p^2\le2p$. Lemmas~\ref{lem:balanced} and~\ref{lem:coding} therefore give
\begin{align*}
 \I(F)&=pm\I(b)+O_{g}(p(s+r)),\\
 \HH(F)&=pm\HH(b)+O_{g}(p(s+r)),\\
 \E F&=1-p,\qquad \Var(F)=2p-p^2.
\end{align*}
It follows that $\I(F)\to0$, $\E F\to1$, and $\HH(F)/\I(F)\to\HH(b)/\I(b)$.

Each unchanged data influence is $p\I_i(b)$, so its ratio to $\Var(F)$ is at most $\I_i(b)/(2-p)\le\eps$. Each new gadget leaf influence is at most $\beta_L\Var(F)\le\eps\Var(F)$. This includes the gadgets replacing the activation controls.

To check spectral distance, partition the final inputs into unchanged singleton blocks and gadget blocks. Two spectral sets coming from different supports of $F_0$ differ on more than $d_0$ blocks, each contributing at least one coordinate to their symmetric difference. Two spectral sets coming from the same outer support must differ within at least one gadget block, where their distance is at least $2t$. Thus the final distance exceeds $2k$.

For $m$ sufficiently large, all the required numerical bounds and the desired ratio hold. The assertion about restrictions follows from Lemma~\ref{lem:plateau} below. Its influence formula can only increase the ratio, while leaving every remaining coordinate influence and the variance unchanged. Deleting at most $k$ coordinates decreases spectral distance by at most $k$, so distance remains greater than $k$. The equality of suprema follows by applying this construction to each seed and letting $\delta$ decrease to zero; the reverse inequality is containment in the full Boolean class.
\end{proof}

\begin{lemma}[A plateau under physical restrictions]\label{lem:plateau}
Suppose $d(\supp\wh F)>q$. For any set $B$ of at most $q$ input coordinates and any assignment $\rho$ to $B$, deleting $B$ is injective on the spectral support and
\begin{equation}\label{eq:plateau}
 \HH(F|\rho)=\HH(F),\qquad
 \I(F|\rho)=\I(F)-\sum_{j\in B}\I_j(F).
\end{equation}
Every remaining coordinate influence is unchanged. If the empty set belongs to the spectral support, then the mean is unchanged as well.
\end{lemma}
\begin{proof}
The coefficient at $T\subseteq B^c$ after restriction is
\[
 \sum_{S:\,S\setminus B=T}\wh F(S)\chi_{S\cap B}(\rho).
\]
Two nonzero summands would correspond to supports at distance at most $|B|\le q$, which is impossible. Thus restriction only changes coefficient signs and relabels their supports injectively. Entropy and the surviving coordinate inclusion probabilities are unchanged. Removing the fixed coordinates gives the influence identity. If the empty support is present, no other support can be contained in $B$, so the empty coefficient is preserved.
\end{proof}

\begin{corollary}[Subexponentially small influence]\label{cor:subexponential}
Let $s:\mathbb N\to[0,\infty)$ satisfy $s(N)=o(N)$. Theorem~\ref{thm:simultaneous} remains valid after adding the requirement
\[
 \I(F)\le2^{-s(N)},
\]
where $N$ is the final number of inputs. In particular, the joint restriction preserves the full optimal FEI constant at the subexponential influence scales of~\cite{Shalev18}.
\end{corollary}
\begin{proof}
Keep the balanced seed on $n$ variables, the gadget with $d_g$ inputs, and the target code distance fixed as in the theorem. Put
\[
 a_m=\max_{1\le D\le2mn}s(D),\qquad
 \ell_m=\lceil a_m+2\log m\rceil.
\]
The assumption on $s$ implies $a_m=o(m)$ and hence $\ell_m=o(m)$; no monotonicity of $s$ is needed. Use $\ell_m$ activation controls, $p_m=2^{-\ell_m}$, and $r_m=O(\log m)$ code coordinates. The final dimension is
\[
 N_m=mn+d_g(\ell_m+r_m)\le2mn
\]
for all sufficiently large $m$. The preceding entropy and influence estimates have errors $O_g(p_m(\ell_m+r_m))=o(p_m m)$, so the seed ratio is preserved. For large $m$, the total influence is at most $p_m m^2\le2^{-a_m}\le2^{-s(N_m)}$. All other bounds in the theorem are unchanged.
\end{proof}

\begin{remark}
For the spectral random set $S_F$, the same distance assumption gives
$\HH(S_F\cap B\mid S_F\setminus B)=0$ for every $|B|\le q$.
This vanishing conditional spectral entropy and the physical restriction plateau are compatible, but they are different statements. A valid induction must also handle later restrictions, after support collisions become possible.
\end{remark}

\section{Multiplexers with shared data variables}\label{sec:multiplexers}

Let $x\in\bits^a$ be a selector and let $(g_b)_{b\in\bits^a}$ be Boolean functions on a common data cube $\bits^n$. Define
\[
 f(x,y)=g_x(y),\qquad N=2^a.
\]
The branches may depend on overlapping data variables. Write $\I_{\rm sel}(f)$ and $\I_{\rm data}(f)$ for the sums of influences in the two groups of inputs.

\begin{proposition}[Shared-data entropy and the collision identity]\label{prop:muxcollision}
For each data Fourier set $T$, put
\[
 R_T=\frac1N\sum_b\wh g_b(T)^2.
\]
When $R_T>0$, define a probability vector and a unit Euclidean vector on the selector labels by
\[
 r_T(b)=\frac{\wh g_b(T)^2}{NR_T},\qquad
 u_T(b)=\frac{\wh g_b(T)}{\sqrt{NR_T}}.
\]
Let $\nu_T$ be the squared-coordinate law of the normalized Walsh--Hadamard transform of $u_T$:
\[
 \nu_T(S)=\left(\frac1{\sqrt N}\sum_b\chi_S(b)u_T(b)\right)^2.
\]
Then
\begin{equation}\label{eq:muxcollision}
 \HH(f)=a+\frac1N\sum_b\HH(g_b)
 +\sum_{T:R_T>0}R_T\bigl[\HH(\nu_T)-\HH(r_T)\bigr],
\end{equation}
where $\HH$ on probability vectors denotes Shannon entropy. In particular,
\begin{equation}\label{eq:muxupper}
 \HH(f)\le2a+\frac1N\sum_b\HH(g_b).
\end{equation}
The influence identities are
\begin{align}
 \I_{\rm data}(f)&=\frac1N\sum_b\I(g_b),\label{eq:muxdata}\\
 \I_{x_j}(f)&=\frac1N\sum_b\P_y\bigl(g_b(y)\ne g_{b^{(j)}}(y)\bigr),\label{eq:muxselector}\\
 \I_{\rm sel}(f)&=\sum_T R_T\sum_S|S|\nu_T(S),\label{eq:muxspectralinfluence}
\end{align}
where $b^{(j)}$ is obtained by flipping the $j$th selector sign.
\end{proposition}
\begin{proof}
The Fourier coefficients of $f$ are
\[
 \wh f(S,T)=\frac1N\sum_b\chi_S(b)\wh g_b(T).
\]
Parseval in the selector coordinates gives $\sum_S\wh f(S,T)^2=R_T$. Thus $R$ is the data-set marginal of the spectral law, and $\nu_T$ is exactly its conditional selector-set law. Consider also the probability law
\[
 \P(B=b,T)=\frac1N\wh g_b(T)^2.
\]
Its entropy is $a+N^{-1}\sum_b\HH(g_b)$, its $T$-marginal is $R$, and its conditional label law is $r_T$. The two applications of the entropy chain rule give
\[
 \HH(R)=a+\frac1N\sum_b\HH(g_b)-\sum_T R_T\HH(r_T),
 \qquad
 \HH(f)=\HH(R)+\sum_T R_T\HH(\nu_T),
\]
which proves~\eqref{eq:muxcollision}. Since $\HH(R)\le a+N^{-1}\sum_b\HH(g_b)$ and $\HH(\nu_T)\le a$, inequality~\eqref{eq:muxupper} follows. Conditioning on the physical selector proves~\eqref{eq:muxdata} and~\eqref{eq:muxselector}; the spectral interpretation gives~\eqref{eq:muxspectralinfluence}.
\end{proof}

\begin{corollary}[Orthogonal neighboring branches]\label{cor:muxorthogonal}
Suppose
\[
 \E[g_b g_{b^{(j)}}]=0\qquad\text{for every }b,j,
 \qquad \HH(g_b)\le C\I(g_b)\qquad\text{for every }b.
\]
Then $\I_{\rm sel}(f)=a/2$ and
\[
 \HH(f)\le\max\{4,C\}\I(f).
\]
More generally, the condition $\I_{\rm sel}(f)\ge ca$ gives the constant $\max\{2/c,C\}$ for any $c>0$.
\end{corollary}
\begin{proof}
For Boolean branches, orthogonality is equivalent to disagreement probability $1/2$. Apply~\eqref{eq:muxselector}, then bound the two terms of~\eqref{eq:muxupper} by $4\I_{\rm sel}(f)$ and $C\I_{\rm data}(f)$, respectively. The last assertion uses $2a\le(2/c)\I_{\rm sel}(f)$.
\end{proof}

\begin{corollary}[Recursive threshold multiplexers]\label{cor:thresholdmux}
Start with regular Boolean functions as leaves. At each later node, select between previously constructed branches using a fresh independent selector, allowing the branches to share data variables. Suppose neighboring branches at each node are orthogonal as in Corollary~\ref{cor:muxorthogonal}. Every resulting function satisfies
\[
 \HH(f)\le13.415\,\I(f).
\]
The constant is independent of the number of branches, the depth, and the pattern of data reuse.
\end{corollary}
\begin{proof}
The companion manuscript~\cite{Threshold} gives the stated constant for every regular leaf. Corollary~\ref{cor:muxorthogonal} preserves it at every node, since $13.415>4$. Induction proves the assertion.
\end{proof}

For a concrete pair with shared data, take
\[
 g_+(y)=\Maj_3(y_1,y_2,y_3),\qquad
 g_-(y)=\Maj_3(y_1,-y_2,y_4).
\]
Their only common nonzero Fourier indices are $\{1\}$ and $\{2\}$, whose inner-product contributions cancel. Thus the two branches are orthogonal despite sharing two inputs. The one-bit multiplexer has two spectral atoms of mass $1/4$ and eight of mass $1/16$, giving $\HH(f)=3$ and $\I(f)=2$.

\begin{remark}
If all branches are identical, $r_T$ is uniform and $\nu_T$ is concentrated at the empty selector set, so~\eqref{eq:muxcollision} cancels the entire label term $a$. For an addressing function, $r_T$ is concentrated at one label and $\nu_T$ is uniform, giving entropy $2a$. A uniform bound for the increment in~\eqref{eq:muxcollision} solely by selector influence is false: Proposition~\ref{prop:bentsplice} supplies Boolean branches for which that ratio is unbounded. Corollary~\ref{cor:muxorthogonal} retains its stated hypothesis.
\end{remark}

\subsection{Nonuniform subcube partitions}

\begin{proposition}[Balanced branches on separate blocks]\label{prop:subcubemux}
Let $(C_b)_b$ be a partition of a selector cube into subcubes of codimensions $d_b$, and put $p_b=2^{-d_b}$. Let $g_b$ be balanced Boolean functions on mutually disjoint data blocks. Set $f(x,y)=g_b(y_b)$ for $x\in C_b$. Then
\begin{align}
 \HH(f)&=2\sum_b p_bd_b+\sum_b p_b\HH(g_b),\label{eq:subcubemuxH}\\
 \I(f)&=\frac12\sum_b p_bd_b+\sum_b p_b\I(g_b).\label{eq:subcubemuxI}
\end{align}
Consequently, a common branch FEI constant $C$ gives the constant $\max\{4,C\}$ for $f$. The partition need not arise from a decision tree. The identities also hold for branches on a common data block when their nonempty Fourier supports are pairwise disjoint.
\end{proposition}
\begin{proof}
The indicator of $C_b$ has $2^{d_b}$ nonzero Fourier coefficients, each of magnitude $p_b$. Every spectral term of the balanced branch $g_b$ contains a nonempty data set, which identifies the branch uniquely under either stated disjointness hypothesis. Its total spectral mass is $p_b$. Conditional on this branch, the selector set is uniform on subsets of the $d_b$ fixed coordinates, independently of the data set with law $P_{g_b}$. The entropy contribution is
\[
 p_b\bigl[\log(1/p_b)+d_b+\HH(g_b)\bigr]
 =p_b\bigl[2d_b+\HH(g_b)\bigr],
\]
and the influence contribution is $p_b[d_b/2+\I(g_b)]$. Summation proves the formulas and the FEI assertion.
\end{proof}

For completeness, unbalanced separate branches have an additional selector-only term. Put
\[
 \mu_b=\E g_b,\quad v_b=1-\mu_b^2,\quad
 m(x)=\sum_b\mu_b1_{C_b}(x),\quad
 \HH^+(g)=\sum_{T\ne\varnothing}\wh g(T)^2\log\frac1{\wh g(T)^2}.
\]
For a real-valued function $u$, write $\HH_{\rm spec}(u)=\sum_T\wh u(T)^2\log(1/\wh u(T)^2)$ and $\I_2(u)=\sum_T|T|\wh u(T)^2$. The same disjoint expansion gives
\begin{align}
 \HH(f)&=\HH_{\rm spec}(m)+\sum_b p_b\bigl[\HH^+(g_b)+2d_bv_b\bigr],\label{eq:unbalancedmuxH}\\
 \I(f)&=\I_2(m)+\sum_b p_b\bigl[\I(g_b)+(d_b/2)v_b\bigr].\label{eq:unbalancedmuxI}
\end{align}
The Fourier expansion of $m$ accounts precisely for those terms with empty data set; all other terms remain separated by their branch blocks.

\section{Conditioned witness forests}\label{sec:forests}

We refine the queue reference of~\cite{Threshold} by retaining whole-forest costs, conditioning on nonempty roots, and subtracting the entropy of interchangeable witnesses. The forest bound itself uses ordered covers of gate supports; the subsequent fractional estimate uses the signed noise bounds for regular gates from~\cite{Threshold}.

\subsection{Ordered covers and a reference law}

Write $f=h_1\vee\cdots\vee h_N$, with a union true when at least one component equals $1$. Remove constant-false gates. A constant-true gate makes $f$ constant, so suppose every remaining gate is nonconstant. Let $R_l$ be its relevant support. Cover $R_l$ by $M_l\ge1$ ordered subsets $R_{l,t}\subseteq R_l$. For $i\in R_{l,t}$ let $r_{l,t}(i)$ be its rank and let $U_i^{l,t}$ be its strict predecessors. Put
\[
 \gamma=\frac65,\quad p_l=\P(h_l=1),\quad
 \pi_l=\min(p_l,1/2),\quad b_l=\log(1/\pi_l).
\]
Choose positive weights satisfying the fractional incidence constraints
\begin{equation}\label{eq:forestpacking}
 \sum_{l:j\in R_l}\lambda_l\le1\qquad(j\in[n]),
\end{equation}
and write
\begin{equation}\label{eq:forestrootconstants}
 L_l=\log(1/\lambda_l),\qquad
 \Lambda=\sum_l\log(1+\lambda_l\pi_l),\qquad
 Z=1-2^{-\Lambda}.
\end{equation}
Each $\lambda_l\le1$. For $s\ge1$ and $i\in R_l$, define
\[
 \theta_s^l(i)=\max_{t:i\in R_{l,t}}
       \min\{1,(s/r_{l,t}(i))^\gamma\},\qquad
 Z_s^l=e_s\bigl((\theta_s^l(i))_{i\in R_l}\bigr),
\]
where $e_s$ is the elementary symmetric polynomial of degree $s$.

\begin{lemma}[Subset normalizer]\label{lem:forestnormalizer}
For $1\le s\le|R_l|$,
\[
 \sum_{i\in R_l}\theta_s^l(i)\le6M_ls,
 \qquad Z_s^l\le(6eM_l)^s.
\]
\end{lemma}
\begin{proof}
For one ordered subset, the sum of its weights is at most
\[
 s+s^\gamma\int_s^\infty t^{-\gamma}\,dt
 =\frac{\gamma}{\gamma-1}s=6s.
\]
Bound the maximum defining $\theta$ by the sum over its $M_l$ orders. Finally $e_s(w)\le(\sum_iw_i)^s/s!$ and $s!\ge(s/e)^s$, proving the claim.
\end{proof}

Define the following random queue.
\begin{enumerate}[label=(\roman*)]
\item Independently activate gate $l$ as a root with probability $\lambda_l\pi_l/(1+\lambda_l\pi_l)$, and queue the roots in a fixed order.
\item At each visit to $l$, choose $s\ge1$ with probability $1/[s(s+1)]$. An invalid size fails. Otherwise choose an $s$-subset $B\subseteq R_l$ with probability $\prod_{i\in B}\theta_s^l(i)/Z_s^l$.
\item Choose $c\ge0$ with probability $1/[(c+1)(c+2)]$, then choose a multiset of $c$ elements of $B$ uniformly. Order its occurrences deterministically. For each occurrence of $j$, select gate $l'$ with probability $\lambda_{l'}$ whenever $j\in R_{l'}$, and append it to the queue. The remaining probability is failure; it is nonnegative by~\eqref{eq:forestpacking}.
\item If the queue empties, output the union of the selected subsets. Failed or infinite runs contribute no output mass.
\end{enumerate}
Let $Q_1(A)$ be the output subprobability. No root is activated with probability $2^{-\Lambda}$. Every finite run with a root outputs a nonempty set. Consequently
\begin{equation}\label{eq:forestconditioning}
 \sum_{A\ne\varnothing}Q_1(A)\le Z,
 \qquad Q_*(A)=Q_1(A)/Z\quad(A\ne\varnothing)
\end{equation}
is a subprobability law. Conditioning on a nonempty initial root set therefore remains valid without any assumption of almost sure termination.

\subsection{Partition forests and their costs}

Fix a nonempty Fourier set $A$ with $P_f(A)>0$. Every $j\in A$ belongs to some $R_l$. Choose an assignment $\tau(j)$ to such a gate, and put
\[
 A_l=\{j\in A:\tau(j)=l\},\quad s_l=|A_l|,
 \quad\mathcal F=\{l:s_l>0\},\quad\kappa_l=|A\cap R_l|.
\]
On $\mathcal F$ draw an edge $l\to l'$ if $l\ne l'$ and $A_l\cap R_{l'}\ne\varnothing$. Choose one root in each source strongly connected component and a directed spanning forest from these roots. Choose a connecting variable in $A_l\cap R_{l'}$ for each parent--child edge. These choices prescribe a finite queue run that visits every gate in $\mathcal F$ once, selects $B_l=A_l$, and outputs $A$. We call such a run a \emph{partition forest}. Its selected subsets partition $A$; arbitrary once-per-gate runs with overlapping selected subsets are not included in this definition.

For $j\in A\cap R_l$, define
\[
 F_l^j(A)=\min_{t:j\in R_{l,t}}
 \ln\frac{r_{l,t}(j)}{|A\cap U_j^{l,t}|+1}.
\]
These quantities are nonnegative. For a partition forest $r$ with root set $\mathcal R$, set
\begin{align}
 B_\lambda(A;r)={}&
 \sum_{l\in\mathcal F}\left[
 L_l+s_l\log M_l+\gamma s_l\log\frac{\kappa_l}{s_l}\right]
 +\sum_{l\in\mathcal R}b_l\notag\\
 &+\frac{\gamma}{\ln2}\sum_{j\in A}F_{\tau(j)}^j(A).\label{eq:forestcost}
\end{align}
Put $B_\lambda(A)=\min_r B_\lambda(A;r)$ and $B_\lambda(\varnothing)=0$. There are finitely many partition forests for each $A$.

\begin{lemma}[Root disjointness]\label{lem:forestrootdisjoint}
The sets $A\cap R_l$ for distinct roots of a partition forest are pairwise disjoint.
\end{lemma}
\begin{proof}
If $j$ lies in two such intersections, its assigned gate $\tau(j)$ has an edge to both roots, except for a possible edge within its own component. A source component cannot receive an edge from a different component. Thus $\tau(j)$ would belong to both source components, a contradiction.
\end{proof}

\begin{theorem}[Conditioned posterior forest bound]\label{thm:posteriorforest}
Put $v=\Var(f)$ and
\[
 K_{\rm ref}=3+2\log3+\log(6e).
\]
For every nonempty spectral $A$, let $\mathcal P_A$ be a probability distribution on distinct actual partition-forest runs producing $A$. Then
\begin{align}
 \HH(f)\le{}&h_2(v)+v(\Lambda+\log Z)+K_{\rm ref}\I(f)\notag\\
 &+\sum_{A\ne\varnothing}P_f(A)
 \left[\E_{r\sim\mathcal P_A}B_\lambda(A;r)
       -\HH(\mathcal P_A)\right].\label{eq:posteriorforest}
\end{align}
In particular,
\begin{equation}\label{eq:minforest}
 \HH(f)\le h_2(v)+v(\Lambda+\log Z)
             +K_{\rm ref}\I(f)+\E B_\lambda(A).
\end{equation}
Here $h_2$ is binary entropy, and the expectation in~\eqref{eq:minforest} is under $P_f$.
\end{theorem}
\begin{proof}
For a specified root set, its selection probability is
\[
 2^{-\Lambda}\prod_{l\in\mathcal R}\lambda_l\pi_l.
\]
Root and child gate names together therefore cost $\sum_{l\in\mathcal F}L_l$: every visited gate is a root or exactly one child. The root rarity cost is $\sum_{l\in\mathcal R}b_l$.

For $i\in A_l$, choose an inner order attaining $F_l^i(A)$. Since $|A\cap U_i^{l,t}|+1\le\kappa_l$,
\[
 \ln^+\frac{r_{l,t}(i)}{s_l}
 \le F_l^i(A)+\ln\frac{\kappa_l}{s_l}.
\]
Lemma~\ref{lem:forestnormalizer} bounds the subset-choice costs by
\[
 |A|\log(6e)+\sum_l s_l\log M_l
 +\gamma\sum_l s_l\log(\kappa_l/s_l)
 +\frac{\gamma}{\ln2}\sum_{j\in A}F_{\tau(j)}^j(A).
\]
The size choices cost at most $|A|\log3$, using $s(s+1)\le3^s$. If $c_l$ is the number of children, then
\[
 |\mathcal R|+\sum_lc_l=|\mathcal F|\le|A|.
\]
The choices of $c_l$ cost at most $|A|(1+\log3)$, since $(c+1)(c+2)\le2\,3^c$. The multiset choices cost at most $2|A|$, using $\binom{s+c-1}{c}\le2^{s+c}$. Thus the probability $q_1(r)$ of this actual run satisfies
\[
 -\log q_1(r)\le\Lambda+K_{\rm ref}|A|+B_\lambda(A;r).
\]
After conditioning on a nonempty root set, $q(r)=q_1(r)/Z$, so its negative logarithm gains $\log Z$.

For a probability law $\mathcal P_A$ on the indicated runs, log-sum gives
\[
 -\log Q_*(A)\le
 \E_{\mathcal P_A}\log\frac1{q(r)}-\HH(\mathcal P_A).
\]
Indeed, normalizing $q$ within this finite collection gives the same inequality with its total mass in place of $Q_*(A)$, and that mass is at most $Q_*(A)$. Count an actual run only once if it has several abstract descriptions. Finally use the subprobability reference with mass $1-v$ at the empty set and mass $vQ_*(A)$ at nonempty $A$. Gibbs' inequality proves~\eqref{eq:posteriorforest}. A point mass at a minimizing forest proves~\eqref{eq:minforest}.
\end{proof}

\begin{lemma}[Weighted activation bound]\label{lem:forestpath}
The quantities above satisfy
\[
 v\Lambda\le(\log e)\I(f).
\]
\end{lemma}
\begin{proof}
Put $p_0=\P(f=-1)$. Independently resample all coordinates of $R_l$ in a fixed order. The initial event $f=-1$ and the final event $h_l=1$ are independent and have joint probability $p_0p_l$. A path between them contains a false-to-true transition. Resampling coordinate $j$ makes such a transition with probability $\I_j(f)/4$, since every intermediate input is uniform. Weight the resulting union bound by $\lambda_l$ and sum:
\[
 p_0\sum_l\lambda_lp_l
 \le\frac14\sum_j\I_j(f)\sum_{l:j\in R_l}\lambda_l
 \le\I(f)/4.
\]
Now $v\le4p_0$, $\pi_l\le p_l$, and $\log(1+x)\le x\log e$ give the result.
\end{proof}

The standard level-entropy estimate, obtained from the edge-isoperimetric inequality, gives $h_2(v)\le2\I(f)$; see~\cite{OWZ11,Threshold}. The conditioning term has a distinct role. For one conjunction of width $n$ with $\lambda=1$, one has
\[
 \Lambda+\log Z=\log\pi=-n,
\]
which cancels its root rarity cost exactly.

\subsection{A fractional bound using pivotal witness distributions}

For this subsection suppose $h_l$ is a Boolean function of $M_l$ regular gates, and take their relevant supports and desirability orders as the ordered cover. A regular gate here means a gate which, after complementing some input coordinates, is monotone and admits a total desirability order: replacing $(x_i,x_j)=(-1,1)$ by $(1,-1)$ cannot decrease its value whenever $i$ precedes $j$. Different gates may use different orders and input orientations.

We use two signed-piece estimates from~\cite{Threshold}. Put $\beta=120/11$. If $d\in\{0,\pm1\}$ is supported on the $j$-pivotal set of one regular inner gate, its normalized Fourier law $\nu$ satisfies
\begin{equation}\label{eq:forestrankinput}
 \nu(F_l^j>u)\le\sqrt{2e}\,e^{-u/2}\qquad(u\ge0).
\end{equation}
If $d$ is supported on the $j$-pivotal set of $h_l$, then
\begin{equation}\label{eq:forestrareinput}
 \nu(\kappa_l<b_l/\beta)
 \le2^{3/5}2^{-5b_l/12}.
\end{equation}
Both statistics are evaluated at $A=T\cup\{j\}$ under the Fourier law of $d$. Arbitrary signs of $d$ are permitted. In~\eqref{eq:forestrareinput}, the density estimate is applied on fibers fixing only coordinates outside $R_l$; it does not assume that the original gate probability survives restrictions inside $R_l$.

For every $j$ with $\alpha_j:=\I_j(f)>0$, split $\D_jf$ pointwise among pivotal outer gates and then among pivotal regular inner gates, with fixed tie rules. Let $d_{j,l,t}$ be the resulting disjoint signed pieces, and put
\[
 q_{j,l,t}=\frac{\|d_{j,l,t}\|_2^2}{\alpha_j},\qquad
 \sum_{l,t}q_{j,l,t}=1,
 \qquad
 S_j(\lambda)=\sum_{l,t}\sqrt{q_{j,l,t}}\,
 M_l^{5/48}\lambda_l^{-11/48}.
\]
Zero pieces are omitted. A pivotal outer gate has a pivotal inner gate, so this decomposition exists.

\begin{theorem}[Fractional pivotal-witness bound]\label{thm:fractionalforest}
With the preceding hypotheses,
\begin{equation}\label{eq:fractionalforest}
 \HH(f)\le C_0\I(f)+\frac{48}{5}
       \sum_j\alpha_j\log S_j(\lambda)+v\log Z,
\end{equation}
where
\begin{align*}
 C_0={}&2+\log e+K_{\rm ref}+\frac{120}{11}\\
 &+\frac{24}{5}\left[\log(\sqrt{2e}+2^{3/5})+\log e\right]
 =40.80487605924\ldots<41.
\end{align*}
\end{theorem}
\begin{proof}
Call $l$ dense at $A$ if $\kappa_l\ge b_l/\beta$, and set
\[
 G_j(l)=\frac{\gamma}{\ln2}F_l^j(A)
       +b_l1_{\{l\text{ not dense}\}},\qquad
 C_j(A)=\min_{l:j\in R_l}\left[G_j(l)+\frac{11}{5}L_l+\log M_l\right].
\]
Choose $\tau$ attaining these minima. Lemma~\ref{lem:forestrootdisjoint} pays the rarity of dense roots by $\beta|A|$. Each nondense root owns at least one coordinate, so its rarity can be charged to the corresponding sum of $G_j$'s. Also
\[
 \sum_l\lambda_l\kappa_l\le|A|.
\]
The log-sum inequality therefore gives
\[
 \sum_l s_l\log(\kappa_l/s_l)\le\sum_l s_lL_l.
\]
Since $\sum_{l\in\mathcal F}L_l\le\sum_ls_lL_l$, the forest cost satisfies
\begin{equation}\label{eq:forestcoarsening}
 B_\lambda(A)\le\beta|A|+\sum_{j\in A}C_j(A).
\end{equation}

Let $\nu_{j,l,t}$ be a piece's normalized Fourier law. Splitting the rank and rarity costs at $u/2$ in~\eqref{eq:forestrankinput}--\eqref{eq:forestrareinput} gives
\[
 \nu_{j,l,t}(G_j(l)>u)\le C_*2^{-\delta u},\qquad
 C_*=\sqrt{2e}+2^{3/5},\quad\delta=5/24.
\]
For $u<0$ the right side exceeds $1$, so the inequality remains valid. For any Fourier event $E$, Minkowski's inequality gives
\[
 \nu_j(E)^{1/2}\le\sum_{l,t}\sqrt{q_{j,l,t}}\,
                                  \nu_{j,l,t}(E)^{1/2},
\]
where $\nu_j(T)=\wh{\D_jf}(T)^2/\alpha_j$. On $E=\{C_j>u\}$, each piece satisfies $G_j(l)>u-(11/5)L_l-\log M_l$. Hence
\[
 \nu_j(C_j>u)\le\min\{1,C_*S_j(\lambda)^2 2^{-\delta u}\}.
\]
The exponents are $\delta/2=5/48$ and $\delta(11/5)/2=11/48$. Since $S_j\ge1$, integration yields
\[
 \E_{\nu_j}C_j\le\frac{24}{5}(\log C_*+\log e)
                       +\frac{48}{5}\log S_j.
\]
Finally $P_f(A)=\alpha_j\nu_j(A\setminus\{j\})$ for $j\in A$. Sum, then apply~\eqref{eq:minforest},~\eqref{eq:forestcoarsening}, Lemma~\ref{lem:forestpath}, and $h_2(v)\le2\I(f)$.
\end{proof}

If the read number is $k$ and $M_l\le M$, choose $\lambda_l=1/k$. There are at most $kM$ nonzero pieces at each coordinate, so $S_j\le k^{35/48}M^{29/48}$. Dropping the nonpositive conditioning term recovers
\[
 \HH(f)\le\left(41+7\log k+\frac{29}{5}\log M\right)\I(f).
\]
Keeping only the outer witness probabilities $q_{j,l}=\sum_tq_{j,l,t}$ gives the intermediate bound
\begin{equation}\label{eq:outerfractional}
 \HH(f)\le\left(41+\frac{29}{5}\log M\right)\I(f)
 +\frac{48}{5}\sum_j\alpha_j
       \log\sum_l\sqrt{q_{j,l}}\lambda_l^{-11/48}.
\end{equation}
For uniform weights, its last term is $(11/5)\I(f)\log k+(24/5)\sum_j\alpha_j\HH_{1/2}(q_j)$, where $\HH_{1/2}(q)=2\log\sum_l\sqrt{q_l}$.

\subsection{Addressing and the cost of repeated labels}

Let $a\ge1$, $N=2^a$, and $f(x,y)=y_x$, with one data bit $y_b$ per selector word $b$. Its canonical DNF has the gates $h_b=[x=b\text{ and }y_b=1]$. Each has probability $1/(2N)$ and support consisting of all selectors and its own data bit. The expansion
\[
 f=\sum_b1_{\{x=b\}}y_b
\]
has $N^2$ coefficients of magnitude $1/N$, giving
\begin{equation}\label{eq:addressingHI}
 \HH(f)=2a,\qquad\I(f)=1+a/2.
\end{equation}
Each selector has influence $1/2$; each data bit has influence $1/N$.

\begin{proposition}[Coordinatewise obstruction and one-root repair]\label{prop:addressforest}
For this DNF, the nonnegative pressure term in~\eqref{eq:outerfractional}, minimized over fractional incidence weights, is exactly
\begin{equation}\label{eq:addresspressure}
 \frac72a^2+\frac{11}{5}a.
\end{equation}
Nevertheless, the conditioned whole-forest bound admits witnesses satisfying
\begin{equation}\label{eq:addressforestpayment}
 \E B_\lambda(A)+\log Z\le a+\frac65\I(f).
\end{equation}
\end{proposition}
\begin{proof}
For a selector derivative, a pivotal input has one true endpoint and one false endpoint. The true endpoint satisfies exactly one gate, so its witness is forced. Each gate's piece has density $1/(2N)$, giving $q_{j,b}=1/N$. A data derivative has a deterministic witness. All selector incidence constraints reduce to $\sum_b\lambda_b\le1$.

For a selector, the sum in~\eqref{eq:outerfractional} is $N^{-1/2}\sum_b\lambda_b^{-11/48}$. Convexity makes it minimal at $\lambda_b=1/N$, with value $N^{35/48}$. The data contribution is minimized at the same weights by concavity of $\log$. The $a$ selectors, each of influence $1/2$, contribute $(7/2)a^2$, and the data coordinates contribute $(11/5)a$. This proves~\eqref{eq:addresspressure}.

For a spectral set $A=\{y_b\}\cup S$, assign every coordinate to gate $b$. There is one root, no children, and $s_b=\kappa_b=|S|+1$, so the allocation cost vanishes. Choose $\lambda_b=1/N$, and order each clause's selectors first in their common order, with its data variable last. The label costs $a$ and the root rarity costs $a+1$.

Conditional on a selector of rank $r$ being selected, its predecessor count is $\operatorname{Bin}(r-1,1/2)$. The exact identity
\[
 \E\frac1{1+\operatorname{Bin}(r-1,1/2)}
 =\frac{2(1-2^{-r})}{r}\le\frac2r
\]
and Jensen's inequality bound its expected $F$-cost by $\ln2$. The same argument with $r=a+1$ applies to the always-present data coordinate. Consequently
\[
 \E\sum_{j\in A}F_b^j(A)\le\I(f)\ln2,
 \qquad \E B_\lambda(A)\le2a+1+\frac65\I(f).
\]
Finally $Z\le\sum_b\lambda_b\pi_b=1/(2N)$, proving~\eqref{eq:addressforestpayment}.
\end{proof}

Thus the quadratic loss in~\eqref{eq:addresspressure} is caused by charging a shared label separately at many Fourier coordinates. It is absent from the whole-forest construction.

\subsection{Posterior label erasure and the remaining condition}

A simple identity illustrates the entropy subtraction in Theorem~\ref{thm:posteriorforest}. Choose uniformly a $w$-subset of $[n]$, and then uniformly an $s$-subset of it, with $s\le w$. The probability of any specified $s$-set $A$ is
\[
 \frac{\binom{n-s}{w-s}}{\binom nw\binom ws}
 =\frac1{\binom ns}.
\]
The posterior root entropy is $\log\binom{n-s}{w-s}$; it subtracts the latent label cost exactly. A single chosen witness would miss this saving.

Theorem~\ref{thm:posteriorforest} leaves the following concrete sufficient condition for unrestricted FEI: a universal finite $K$ and, for every Boolean $f$, ordered gate representations, fractional weights, and realizable partition-forest distributions satisfying
\begin{equation}\label{eq:foresttarget}
 \boxed{\quad
 v(\Lambda+\log Z)+
 \sum_{A\ne\varnothing}P_f(A)
 \left[\E_{r\sim\mathcal P_A}B_\lambda(A;r)-\HH(\mathcal P_A)\right]
 \le K\I(f).
 \quad}
\end{equation}
This would give the FEI constant $2+K_{\rm ref}+K$. The condition remains open. If all successful runs are permitted, successively larger finite collections identify the variational infimum with the cross-entropy of the output reference. This does not bound that cross-entropy. Nor does a generic minterm representation suffice: the uniform reference on nonempty subsets gives only
\[
 \HH(f)\le h_2(v)+v\log(2^n-1)\le2\I(f)+nv.
\]
For balanced odd majority, $v=1$ while $\I(f)\le\sqrt n$, so this dimension term cannot be absorbed universally. An estimate such as~\eqref{eq:foresttarget} must account for the sharing and erasure of witnesses using additional structure of Boolean spectra.

\section{Restriction increments and coordinate selection}\label{sec:restrictions}

For a coordinate $i$, let $f_i^+$ and $f_i^-$ denote the two restrictions
of $f$, and put
\[
 \Delta_i\HH(f)=\HH(f)-\frac{\HH(f_i^+)+\HH(f_i^-)}2.
\]
The restriction framework of Gonz\'alez, MacManus and
Pereyra~\cite{GMP26} asks for a coordinate at which this entropy
increment can be charged to the influence lost in the split.
We first give an exact scalar description of the increment, then
identify several limitations of possible coordinate rules.

\subsection{The exact increment and its scalar profile}

Write $h_2(u)=-u\log u-(1-u)\log(1-u)$, with its continuous endpoint
values, and define
\begin{equation}\label{eq:kappadef}
 \kappa(u)=1+h_2(u)
 -h_2\!\left(\frac12+\sqrt{u(1-u)}\right),
 \qquad 0\le u\le1.
\end{equation}
For $S\subseteq[n]\setminus\{i\}$ set
\[
 c_S=\wh f(S),\qquad d_S=\wh f(S\cup\{i\}),\qquad
 m_S=c_S^2+d_S^2,\qquad u_S=\frac{d_S^2}{m_S},
\]
omitting terms with $m_S=0$.

\begin{proposition}[Exact restriction increment]\label{prop:splitexact}
For every Boolean $f$ and every coordinate $i$,
\begin{equation}\label{eq:splitexact}
 \Delta_i\HH(f)=\sum_{S\not\ni i}m_S\kappa(u_S).
\end{equation}
In particular, $0\le\Delta_i\HH(f)\le2$.

If $\mathcal S$ has law $P_f$, define
\[
 B_i(f)=H\bigl(\mathbf1_{\{i\in\mathcal S\}}\mid
                         \mathcal S\setminus\{i\}\bigr).
\]
Then $\Delta_i\HH(f)=B_i(f)+M_i(f)$, where
\begin{equation}\label{eq:mixingbound}
 0\le M_i(f)\le
 \frac2{\ln2}\sum_{S\not\ni i}\min(c_S^2,d_S^2)
 \le\frac2{\ln2}\I_i(f).
\end{equation}
\end{proposition}

\begin{proof}
The two restricted coefficients indexed by $S$ are $c_S+d_S$ and
$c_S-d_S$. Suppressing subscripts and putting
$v=(c+d)^2/(2m)=1/2+cd/m$, the parent pair contributes
$-m\log m+mh_2(u)$ to entropy, whereas the average contribution
of the two branches is $-m\log(2m)+mh_2(v)$.
Their difference is $m[1+h_2(u)-h_2(v)]$.
Changing the sign of $cd$ does not affect $h_2(v)$, proving
\eqref{eq:splitexact}. Both $h_2(u)$ and $1-h_2(v)$ lie in $[0,1]$,
and $\sum_Sm_S=1$.

The term $\sum_Sm_Sh_2(u_S)$ is exactly $B_i(f)$, leaving
\[
 M_i(f)=\sum_Sm_S
 \left[1-h_2\!\left(\frac12+\frac{c_Sd_S}{m_S}\right)\right].
\]
To bound it, let
\[
 g(t)=\frac{(1+t)\ln(1+t)+(1-t)\ln(1-t)}2,
 \qquad 0\le t\le1.
\]
For $0\le t<1$, the inequality
$\atanh t\le t/\sqrt{1-t^2}$ follows by comparing derivatives
from zero. Since $g'(t)=\atanh t$, integration gives
$g(t)\le1-\sqrt{1-t^2}$. Apply this with
$t=2\sqrt{u(1-u)}$, for which
$1-\sqrt{1-t^2}=2\min(u,1-u)$, and divide by $\ln2$.
Continuity handles the endpoints. Multiplication by $m_S$ and
summation yield \eqref{eq:mixingbound}.
\end{proof}

\begin{lemma}[Concavity of the scalar profile]\label{lem:kappaconcavity}
The function $\kappa$ is continuous on $[0,1]$, symmetric about
$1/2$, and strictly concave on $(0,1)$. It satisfies
$\kappa(0)=\kappa(1)=0$ and $\kappa(1/2)=2$.
Consequently,
\begin{equation}\label{eq:kappabound}
 0\le\Delta_i\HH(f)\le\kappa(\I_i(f)).
\end{equation}
As $u\downarrow0$,
\begin{equation}\label{eq:kappasmall}
 \kappa(u)=u\log(1/u)+\frac{3u}{\ln2}+O(u^2).
\end{equation}
\end{lemma}

\begin{proof}
Put $t=2\sqrt{u(1-u)}$ and
$G(t)=1-h_2((1+t)/2)$. Away from the midpoint,
\[
 h_2''(u)=-\frac4{(\ln2)t^2},\qquad
 G'(t)=\frac{\atanh t}{\ln2},\qquad
 G''(t)=\frac1{(\ln2)(1-t^2)},
\]
\[
 t'=\frac{2(1-2u)}t,\qquad t''=-\frac4{t^3}.
\]
The contributions $h_2''(u)$ and $G''(t)(t')^2$ cancel, leaving
\begin{equation}\label{eq:kappasecond}
 \kappa''(u)=-
 \frac{4\atanh(2\sqrt{u(1-u)})}
 {(\ln2)(2\sqrt{u(1-u)})^3}<0.
\end{equation}
At the midpoint, with $u=1/2+x$, the expansion is
\[
 \kappa(1/2+x)
 =2-x^2\log(1/x^2)-\frac{3x^2}{\ln2}
       +O\bigl(x^4\log(1/|x|)\bigr).
\]
Thus the first derivative tends to zero on both sides; the
concavity on the two half-intervals joins across $1/2$.
Continuity and the endpoint values follow from
\eqref{eq:kappadef}. Jensen's inequality, together with
$\sum_Sm_S=1$ and $\sum_Sm_Su_S=\I_i(f)$, proves
\eqref{eq:kappabound}. Finally,
$h_2(u)=u\log(1/u)+u/\ln2+O(u^2)$ and
$G(2\sqrt{u(1-u)})=2u/\ln2+O(u^2)$ give
\eqref{eq:kappasmall}.
\end{proof}

Since a concave function vanishing at zero has a nonincreasing
ratio to its argument, any coordinate with $\I_i(f)\ge\tau>0$
satisfies
\begin{equation}\label{eq:kappacutoff}
 \Delta_i\HH(f)\le\frac{\kappa(\tau)}{\tau}\I_i(f).
\end{equation}
For example, $\kappa(1/16)/(1/16)=8.21872678758653\ldots$.
The logarithm in \eqref{eq:kappasmall} cannot be replaced by a
uniform constant for every coordinate, as
Proposition~\ref{prop:bentsplice} below shows.

\subsection{Valid induction and block exposure}

The exact influence identity is
\begin{equation}\label{eq:restrictionI}
 \I(f)=\I_i(f)+\frac{\I(f_i^+)+\I(f_i^-)}2.
\end{equation}
Indeed, restriction averages preserve the influence of every
other coordinate, while coordinate $i$ disappears.

\begin{proposition}[A suitable-coordinate induction]\label{prop:restrictioninduction}
Suppose a class of Boolean functions is closed under restrictions
and each nonconstant member has a relevant coordinate $i$ with
$\Delta_i\HH(f)\le C\I_i(f)$. Then every member satisfies
$\HH(f)\le C\I(f)$.

Existence of a finite global FEI constant is also equivalent to
existence of a finite $K$ such that every nonconstant Boolean
function has a relevant coordinate satisfying $B_i(f)\le K\I_i(f)$.
A global FEI constant $C$ permits $K=C$; the latter erasure
inequality implies FEI with constant $K+2/\ln2$.
\end{proposition}

\begin{proof}
The first assertion follows by induction on the number of
coordinates, using \eqref{eq:restrictionI}.
For the second, the entropy chain rule gives
\begin{equation}\label{eq:hanerasure}
 \sum_i B_i(f)\le\HH(f),
\end{equation}
because conditioning on all other coordinates can only decrease
each summand of a fixed chain rule. Under $\HH(f)\le C\I(f)$,
some positive-influence coordinate must have $B_i(f)\le C\I_i(f)$.
Conversely, \eqref{eq:mixingbound} turns the proposed erasure
charge into a suitable-coordinate entropy charge, to which the
first assertion applies.
\end{proof}

The quantifier is existential at every restricted function.
Neither an arbitrary coordinate nor a preferred coordinate rule
is justified by Proposition~\ref{prop:restrictioninduction}.
The implication is hereditary: it is not the pointwise
inequality
\[
 \HH(f)\le\sum_iB_i(f)+\frac2{\ln2}\I(f),
\]
which is false for the coded examples in~\cite{BrainErasure}.
Theorem~\ref{thm:simultaneous} and Lemma~\ref{lem:plateau}
also show how finite depths of restrictions can preserve entropy
while revealing no local entropy cost.

More generally, consider a finite adaptive restriction tree, and
choose its branches by the uniform input. If $f_v$ is the function
at a visited internal node and $i(v)$ its split coordinate, then
\begin{align}
 \HH(f)&=\E\sum_{v\ {\rm internal\ on\ the\ path}}
               \Delta_{i(v)}\HH(f_v)
              +\E\HH(f_{\rm leaf}),\label{eq:treetelescopeH}\\
 \I(f)&=\E\sum_{v\ {\rm internal\ on\ the\ path}}
               \I_{i(v)}(f_v)
              +\E\I(f_{\rm leaf}).\label{eq:treetelescopeI}
\end{align}
These identities permit amortized charges and stopping at classes
with established FEI bounds, as in~\cite{GMP26}. They require the
cost along the whole path to be controlled.

\begin{proposition}[Exact block increment]\label{prop:blockincrement}
Fix a block $W\subseteq[n]$. For $T\subseteq W^c$ put
\[
 q_T=\sum_{U\subseteq W}\wh f(T\cup U)^2,\qquad
 \phi_T(y)=q_T^{-1/2}
       \sum_{U\subseteq W}\wh f(T\cup U)\chi_U(y),
\]
omitting $q_T=0$. Then
\begin{equation}\label{eq:blockincrement}
 \HH(f)-\E_y\HH(f|_{W=y})
 =H(\mathcal S\cap W\mid\mathcal S\cap W^c)
   +\sum_Tq_T\operatorname{Ent}_2(\phi_T^2),
\end{equation}
where
$\operatorname{Ent}_2 Z=\E[Z\log(Z/\E Z)]$. Moreover,
\begin{equation}\label{eq:blockmixing}
 0\le\sum_Tq_T\operatorname{Ent}_2(\phi_T^2)
 \le\frac2{\ln2}\sum_{i\in W}\I_i(f).
\end{equation}
\end{proposition}

\begin{proof}
Let $p_T(U)=\wh f(T\cup U)^2/q_T$. Then
$\E\phi_T^2=1$, and the original entropy is
$H(q)+\sum_Tq_TH(p_T)$. The restricted coefficient indexed by
$T$ is $\sqrt{q_T}\phi_T(y)$, so its entropy averaged over $y$
is
\[
 H(q)-\sum_Tq_T\operatorname{Ent}_2(\phi_T^2).
\]
Subtracting proves \eqref{eq:blockincrement}. The uniform-cube
logarithmic Sobolev inequality~\cite{OD14}, in the normalization
$D_jg=(g|_{x_j=1}-g|_{x_j=-1})/2$, states
\[
 \operatorname{Ent}_2(g^2)
 \le\frac2{\ln2}\sum_j\E(D_jg)^2.
\]
Applying it to each $\phi_T$ and using Parseval gives
\[
 \sum_Tq_T\sum_{i\in W}\E(D_i\phi_T)^2
 =\sum_{T,U}|U|\wh f(T\cup U)^2
 =\sum_{i\in W}\I_i(f),
\]
which proves \eqref{eq:blockmixing}.
\end{proof}

Block exposure therefore retains a conditional erasure charge.
Taking $W=[n]$ makes that charge $\HH(f)$ itself. Small erasure
for bounded blocks at the initial function does not replace the
hereditary requirement, because the separated supports in
Lemma~\ref{lem:plateau} have zero entropy increment under all
sufficiently shallow restrictions.

\subsection{Character-centered budgets}

Define
\[
 \D(f)=\sum_j t(\I_j(f)),\qquad t(u)=\min(u,1-u).
\]
Character multiplication preserves entropy and complements the
influences on its selected coordinates. Hence
$\D(f)=\min_U\I(\chi_Uf)$ and $\D$ is character invariant.
For a split define
\[
 d_i^{\D}(f)=\D(f)
            -\frac{\D(f_i^+)+\D(f_i^-)}2.
\]

\begin{proposition}[Centered budget and its collapse]\label{prop:centerbudget}
For every $f,i$,
\begin{align}
 d_i^{\D}(f)
 ={}&t(\I_i(f))\notag\\
 &+\sum_{j\ne i}
 \left[t(\I_j(f))
       -\frac{t(\I_j(f_i^+))+t(\I_j(f_i^-))}2\right]
 \ge t(\I_i(f)).\label{eq:centerbudget}
\end{align}
Both $d_i^{\D}$ and $\Delta_i\HH$ are character invariant.
If $\max_j\I_j(f)\le1/4$, then
\begin{equation}\label{eq:centerbudgetcollapse}
 d_i^{\D}(f)=\I_i(f)\qquad\text{for every }i.
\end{equation}
\end{proposition}

\begin{proof}
For $j\ne i$, the two child influences average to $\I_j(f)$.
Concavity of $t$ proves \eqref{eq:centerbudget}.
Restriction commutes with character multiplication up to a
character and an output sign, proving the invariances.
Each child influence is at most twice its parent's influence.
Under the stated upper bound, all parent and child influences
are at most $1/2$, so $\D=\I$ at these three nodes.
Equation~\eqref{eq:restrictionI} now gives
\eqref{eq:centerbudgetcollapse}.
\end{proof}

A suitable-coordinate inequality
$\Delta_i\HH\le C d_i^{\D}$ would telescope to
$\HH\le C\D$, stopping at characters, whose entropy is zero.
The centered budget is sometimes larger than the original
budget: for a two-input conjunction it reduces the local ratio
from four to two. It supplies no improvement in
\eqref{eq:centerbudgetcollapse}. More generally, if
$\max_j\I_j(f)\le2^{-(r+1)}$, every restriction fixing at most
$r$ coordinates has influences at most $1/2$, so the first
$r$ levels have no centering gain. This follows by iterating
the factor-two bound. The low-influence reductions
in~\cite{BrainCritical} and Theorem~\ref{thm:simultaneous}
therefore remain applicable.

\subsection{Bad coordinates with arbitrarily small influences}

The following bent-splice mechanism also underlies the local
allocation obstruction in~\cite{BrainAllocation}.

\begin{proposition}[Bent splice]\label{prop:bentsplice}
For each even $m\ge4$, there is a Boolean function $F$ on $m+1$
coordinates, with distinguished coordinate $z$, such that
\[
 \I_z(F)=2^{-m},\qquad
 \I_j(F)=\frac12+2^{-m}\quad(1\le j\le m),
\]
and
\begin{equation}\label{eq:bentspliceratio}
 \frac{\Delta_z\HH(F)}{\I_z(F)}
 \ge\log\!\left(2^m-2\cdot2^{m/2}+2\right)\ge m-1.
\end{equation}
\end{proposition}

\begin{proof}
Use binary coordinates for the $m$ base variables, with characters
$(-1)^{S\cdot x}$, and put
\[
 g(x)=(-1)^{x_1x_2+x_3x_4+\cdots+x_{m-1}x_m},\qquad
 h(x)=g(x)-2\mathbf1_{\{x=0\}}.
\]
Both are Boolean. Define $F(x,1)=g(x)$ and $F(x,-1)=h(x)$.
Set $N=2^m$, $\alpha=N^{-1}$ and $R=\sqrt N$.
The two-bit quadratic factor has four Fourier coefficients of
magnitude $1/2$; tensorization gives
$\wh g(S)\in\{-1/R,1/R\}$ for every $S$.
Each base influence of $g$ is $1/2$.
Since $g(0)=g(e_j)=1$, flipping the value at zero changes the
edge $\{0,e_j\}$ from agreeing to disagreeing, increasing
$\I_j(h)$ by $2/N$. Averaging the branches proves the displayed
influence identities.

For the $z$-split, the coefficient pairs are
$c_S=\wh g(S)-\alpha$ and $d_S=\alpha$. Thus
\[
 \frac{m_S}{d_S^2}=(\pm R-1)^2+1
                   \ge N-2\sqrt N+2.
\]
The $d_S^2\log(m_S/d_S^2)$ term alone gives
\[
 \Delta_z\HH(F)\ge B_z(F)
 \ge N\alpha^2\log(N-2\sqrt N+2).
\]
Since $\I_z(F)=\alpha$, the first inequality in
\eqref{eq:bentspliceratio} follows. The second follows from
$N-2\sqrt N+2-N/2=(\sqrt N-2)^2/2\ge0$.
\end{proof}

\begin{corollary}[Failure of two maximum-budget rules]\label{cor:badmaximum}
There are Boolean functions $\widetilde F_m$ with
$\max_i\I_i(\widetilde F_m)\le2^{-m}$ for which the unique
maximum-influence coordinate has local ratio at least $m-1$.
It is also the unique coordinate maximizing $d_i^{\D}$, with
the same ratio $\Delta_i\HH/d_i^{\D}$.
\end{corollary}

\begin{proof}
In Proposition~\ref{prop:bentsplice}, replace each base coordinate
by a majority on a disjoint block of
$K=2N^2+1$ bits, keeping $z$ unchanged. Its leaf influence is
\[
 \beta_K=\binom{K-1}{(K-1)/2}2^{-(K-1)}
 \le\frac1{\sqrt{(K+1)/2}}<\alpha.
\]
For completeness, the bound
$\binom{2t}{t}/4^t\le1/\sqrt{t+1}$ follows by induction:
the consecutive ratio is $(2t+1)/(2t+2)$, and the squared
comparison reduces to $3t+2\ge0$.

Lemma~\ref{lem:balanced} gives each new leaf influence
$(1/2+\alpha)\beta_K<\alpha$ and preserves $\I_z=\alpha$.
It also preserves $\Delta_z\HH$: the additional entropy
in the parent is $\sum_j\I_j(F)\HH(\Maj_K)$, and the
additional branch entropies average to the same quantity
by \eqref{eq:restrictionI} applied to the other coordinates.
Thus $z$ is the unique maximum-influence coordinate, and
\eqref{eq:bentspliceratio} persists.
All influences are at most $\alpha\le1/16$, so
Proposition~\ref{prop:centerbudget} gives
$d_i^{\D}=\I_i$ at every coordinate, proving the second assertion.
\end{proof}

These examples reject the two specified rules, not the
existential hypothesis of Proposition~\ref{prop:restrictioninduction}.
The cost may still be paid by choosing a different coordinate
or by an amortized argument.

\subsection{Separable influence potentials}

\begin{proposition}[An obstruction to charging every split]\label{prop:separablepotential}
Let $\phi:[0,1]\to\mathbb R$ satisfy $\phi(0)=0$,
$\phi(t)=O(t)$ as $t\downarrow0$, and have a finite left derivative
at $1/2$. Put $P_\phi(f)=\sum_j\phi(\I_j(f))$.
There is no finite $C>0$ such that
\begin{equation}\label{eq:falsepotentialcharge}
 \Delta_i\HH(f)\le C\left[
 P_\phi(f)-\frac{P_\phi(f_i^+)+P_\phi(f_i^-)}2\right]
\end{equation}
for every Boolean $f$ and every coordinate $i$.
The assertion remains false after adding to $P_\phi(f)$ any
real-valued function of $\E f$.
\end{proposition}

\begin{proof}
Before flipping zero in Proposition~\ref{prop:bentsplice},
multiply $g$ by the character on all its base coordinates.
Its spectrum remains flat, and its base influences remain
$1/2$. Now the edges from zero to $e_j$ disagree, so flipping
zero makes the other branch's influences $1/2-2\alpha$ and
the parent influences $1/2-\alpha$.
The selector still has influence $\alpha$, and the lower bound
$\Delta_z\HH\ge(m-1)\alpha$ is unchanged.
The potential drop is exactly
\[
 \phi(\alpha)+m\left[
 \phi(1/2-\alpha)
 -\frac{\phi(1/2)+\phi(1/2-2\alpha)}2\right].
\]
The bracket is $o(\alpha)$ by the one-sided derivative
assumption. Dividing the displayed expression by $m\alpha$
therefore gives a quantity tending to zero, whereas
$\Delta_z\HH/(m\alpha)\ge1-1/m$.
This contradicts \eqref{eq:falsepotentialcharge}.

Finally multiply the entire splice by a fresh independent sign.
Parent and both selector branches then have mean zero, so any
mean-only correction cancels. The new coordinate has influence
one at all three nodes, and its contribution to the potential
drop also cancels. Entropy and the selector increment are
unchanged.
\end{proof}

The assumptions include $\phi(t)=t$ and
$\phi(t)=\min(t,1-t)$. The proposition concerns a charge at
every coordinate; it does not rule out a suitable-coordinate
inequality, a nonseparable potential, or a global tree budget.

\subsection{Finite checks and exact witnesses}\label{subsec:restrictionchecks}

Enumeration of all $2^{16}=65,536$ Boolean truth tables on four
labelled coordinates, including functions with unused inputs,
gives the values in Table~\ref{tab:restrictionchecks}.
Walsh numerators, influences and centered budgets are calculated
with exact integer or dyadic arithmetic. The ordering of
entropy ratios in the enumeration uses floating-point
logarithms; the exact witness formulas below were also evaluated
with 80-digit arithmetic. Thus the table reports numerical
finite maxima, not interval-certified comparisons or
dimension-independent bounds.

\begin{table}[ht]
\centering
\caption{Observed maxima among four-variable Boolean functions.
Ratios with zero denominator are omitted.}\label{tab:restrictionchecks}
\begin{tabular}{lc}
\toprule
Quantity & Observed maximum\\
\midrule
$\HH(f)/\I(f)$ & $3.402475551198587$\\
$\min_{i:\I_i>0}\Delta_i\HH(f)/\I_i(f)$ & $4.867668746754797$\\
$\min_{i:d_i^{\D}>0}\Delta_i\HH(f)/d_i^{\D}(f)$
 & $4.867668746754797$\\
$\max_{i:\I_i>0}\Delta_i\HH(f)/\I_i(f)$
 & $5.921783780331384$\\
$\max_{i:d_i^{\D}>0}\Delta_i\HH(f)/d_i^{\D}(f)$
 & $5.921783780331384$\\
\bottomrule
\end{tabular}
\end{table}

The entropy-to-influence witness is a four-input conjunction,
up to output sign, with exact ratio
\[
 12-\frac{49}{16}\log 7.
\]
The suitable-coordinate witness is a three-input conjunction
and one unused variable. Its squared coefficients are $9/16$
at the empty set and $1/16$ at each of seven nonempty sets;
each relevant influence and centered budget is $1/4$.
Its local ratio is exactly
\[
 12-\frac92\log 3
 =4.86766874675479718345817475223482571058\ldots .
\]
For an arbitrary-coordinate witness, take the function on
$\{0,1\}^4$ equal to $-1$ precisely at
$0000,0001,0010$. Its squared coefficients are $25/64$ once,
$9/64$ three times, and $1/64$ twelve times. At either of its
coordinates of influence $1/8$, the local ratio is
\[
 24-\frac{25}{4}\log 5-\frac94\log 3
 =5.92178378033138391753959094180872425614\ldots .
\]
The exact formulas establish these witness values independently
of the numerical ordering. Larger self-composition examples
in~\cite{BrainCritical} require larger suitable-coordinate
constants, so the finite table cannot be used as a uniform
restriction bound.

\section{Product measures and obstructions to threshold normalization}
\label{biasfull:section}

The generalized FEI inequality of O'Donnell and Tan~\cite{ODT13}
uses a corrected spectral entropy under product measures. Symmetric
threshold gates satisfy this inequality under arbitrary product inputs
in~\cite{Threshold}. The normalization estimates used for those gates
cannot be extended to weighted thresholds without a different argument.
We give explicit counterexamples to both a conditional-total-correlation
estimate and its proposed pivotal-credit replacement. Their genuine
corrected entropy-to-deficit ratios tend to \(2\).

\subsection{Normalization and exact composition}

Let \(Y_i\in\{0,1\}\) be independent with
\(\P(Y_i=1)=q_i\in(0,1)\), and let
\[
 \phi_i(Y_i)=\frac{Y_i-q_i}{\sqrt{q_i(1-q_i)}},
 \qquad v_i=4q_i(1-q_i),\qquad v_S=\prod_{i\in S}v_i.
\]
For a Boolean sign function \(f\), put
\[
 P_f^\mu(S)=\wh f_\mu(S)^2,\qquad
 \I_\mu(f)=\sum_S|S|P_f^\mu(S),\qquad
 V_\mu(f)=\Var_\mu(f),
 \qquad \Delta_\mu(f)=\I_\mu(f)-V_\mu(f).
\]
Its corrected entropy is
\begin{equation}\label{biasfull:hplus}
 \HH_\mu^+(f)=\sum_{S\ne\varnothing}P_f^\mu(S)
                      \log\frac{v_S}{P_f^\mu(S)}.
\end{equation}
Every summand is nonnegative, since
\[
 |\wh f_\mu(S)|\le\E_\mu|\phi_S|=\sqrt{v_S}.
\]
Write
\[
 h_0(f)=(\E_\mu f)^2\log\frac1{(\E_\mu f)^2}.
\]
If \(\HH_\mu(f)\) denotes the ordinary Shannon entropy of
\(P_f^\mu\), then
\begin{equation}\label{biasfull:rawcorrected}
 \HH_\mu^+(f)=\HH_\mu(f)-h_0(f)
                    +\sum_i\I_i^\mu(f)\log v_i.
\end{equation}
In particular, on a uniform cube
\(\HH^+(f)=\HH(f)-h_0(f)\).

We use the following exact form of disjoint composition.

\begin{lemma}[Product-measure composition]\label{biasfull:composition}
Let \(g_i\) be Boolean functions on disjoint variable blocks with product
input laws \(\mu_i\). Let \(p_i=\P(g_i=1)\), \(V_i=4p_i(1-p_i)>0\), and
let \(\nu\) be the product law with these output probabilities.
For \(F=G(g_1,\ldots,g_k)\), set \(a_i=\I_i^\nu(G)/V_i\).
Then
\begin{align}
 \HH_\mu^+(F)&=\HH_\nu^+(G)+\sum_i a_i\HH_{\mu_i}^+(g_i),
       \label{biasfull:comph}\\
 \I_\mu(F)&=\sum_i a_i\I_{\mu_i}(g_i),
       \label{biasfull:compi}\\
 \Delta_\mu(F)&=\Delta_\nu(G)+\sum_i a_i\Delta_{\mu_i}(g_i).
       \label{biasfull:compdelta}
\end{align}
Constant children may be frozen and deleted.
\end{lemma}
\begin{proof}
For a nonempty spectral set, its occupied blocks form a unique set \(A\),
and its nonempty intersections with those blocks are \(T_i\), \(i\in A\).
The coefficient is
\[
 \wh F\Bigl(\bigcup_{i\in A}T_i\Bigr)
   =\wh G_\nu(A)\prod_{i\in A}\frac{\wh g_i(T_i)}{\sqrt{V_i}}.
\]
The squared coefficient is therefore the outer spectral mass times the
conditional nonempty spectral masses of the occupied children.
Taking the logarithm in~\eqref{biasfull:hplus} proves
\eqref{biasfull:comph}; summing degrees proves~\eqref{biasfull:compi}.
Finally \(V_\mu(F)=V_\nu(G)\), so subtraction proves
\eqref{biasfull:compdelta}.
\end{proof}

For an increasing event \(A\), write \(f=2\mathbf1_A-1\),
\(\alpha=\mu(A)\), and let \(\rho_i\) be its pivotal probability.
Then
\begin{equation}\label{biasfull:pivotal}
 \I_i^\mu(f)=4q_i(1-q_i)\rho_i,\qquad
 \E_{\mu(\cdot\mid A)}Y_i-q_i
       =\frac{\I_i^\mu(f)}{4\alpha}.
\end{equation}
Both identities follow by conditioning on \(Y_i\).
For a distribution \(Q\) on a finite product set, its total correlation,
in nats, is
\[
 \TC(Q)=\KL\left(Q\,\middle\|\,\bigotimes_iQ_i\right)
       =\sum_i H_{\rm nat}(Q_i)-H_{\rm nat}(Q).
\]
If \(Q=\mu(\cdot\mid A)\), the product structure of \(\mu\) gives
\begin{equation}\label{biasfull:tcdecomp}
 \ln\frac1\alpha
   =\TC(Q)+\sum_i\KL(Q_i\|\mu_i).
\end{equation}

\subsection{Rare disjoint tribes}

\begin{proposition}\label{biasfull:tribes}
Let
\[
 A_m=\bigvee_{j=1}^m(X_j\wedge Z_j),\qquad f_m=2\mathbf1_{A_m}-1,
\]
where the \(2m\) inputs are independent Bernoulli \(q=m^{-\gamma}\)
bits and \(1/2<\gamma<1\).
Writing \(\alpha=\P(A_m)\), \(V=V_\mu(f_m)\), and \(Q=\mu(\cdot\mid A_m)\),
one has
\begin{align}
 \frac{\alpha\TC(Q)}{\I_\mu(f_m)}
     &\sim\frac{2\gamma-1}{8}\ln m,\label{biasfull:tribestc}\\
 \frac{V\ln(1/V)-\sum_i\I_i^\mu(f_m)\ln^+(4\rho_i/V)}
           {\I_\mu(f_m)}
     &\sim\frac{2\gamma-1}{2}\ln m,\label{biasfull:tribescredit}\\
 \frac{\HH_\mu^+(f_m)}{\Delta_\mu(f_m)}
     &\longrightarrow2.\label{biasfull:tribesratio}
\end{align}
For all sufficiently large \(m\), every logarithmic credit in
\eqref{biasfull:tribescredit} is zero.
\end{proposition}
\begin{proof}
Put \(b=q^2\), \(a=(1-b)^m\), \(r=(1-b)^{m-1}\), and
\(\alpha=1-a\). Each pivotal probability is \(qr\), whence
\begin{equation}\label{biasfull:tribesIV}
 \I_\mu(f_m)=8mq^2(1-q)r,\qquad V=4\alpha a.
\end{equation}
Since \(mq^2\to0\) and \(mq\to\infty\),
\[
 \alpha\sim mq^2,\qquad
 \I_\mu(f_m)\sim8\alpha,\qquad
 V\sim4\alpha,\qquad \Delta_\mu(f_m)\sim4\alpha.
\]
The logarithm arguments are exactly
\[
 \frac{4\rho_i}{V}=\frac{q}{\alpha(1-q^2)}
       \sim\frac1{mq}\longrightarrow0.
\]
This proves~\eqref{biasfull:tribescredit} and its final assertion.

Every conditional marginal under \(Q\) has mean \(q+d\), where
\[
 d=\frac{q^2(1-q)(1-q^2)^{m-1}}\alpha\sim\frac1m.
\]
Equation~\eqref{biasfull:tcdecomp} gives
\[
 \TC(Q)=\ln(1/\alpha)
       -2m\KL\bigl(\operatorname{Ber}(q+d)\|
                         \operatorname{Ber}(q)\bigr).
\]
The inequality \(\ln t\le t-1\) implies
\[
 \KL\bigl(\operatorname{Ber}(q+d)\|
                         \operatorname{Ber}(q)\bigr)
       \le\frac{d^2}{q(1-q)}.
\]
Thus the subtracted term is \(O((mq)^{-1})=o(1)\).
This proves~\eqref{biasfull:tribestc}.

For~\eqref{biasfull:tribesratio}, the sign version \(g\) of a single
two-input conjunction has
\begin{align}
 \Delta_\mu(g)&=4q^2(1-q)^2,\notag\\
 \HH_\mu^+(g)&=8q^2(1-q)^2
                  +16q^3(1-q)\log(1/q).
                    \label{biasfull:andpair}
\end{align}
These follow by listing its two singleton masses
\(4q^3(1-q)\) and its two-coordinate mass \(4q^2(1-q)^2\).
The outer OR gate has \(m\) inputs of probability \(b\).
Its variance and influence are \(V_o=4a\alpha\) and
\(\I_o=4mba\), and direct summation of its spectrum gives
\begin{equation}\label{biasfull:orentropy}
 \HH_o^+=2(\I_o-V_o)
           +8ma(\alpha-b)\log\frac1{1-b}.
\end{equation}
Indeed its nonempty spectral mass at a set of size \(k\) is
\(4a^2[b/(1-b)]^k\).
The union bound and the first two Bonferroni terms give
\[
 0\le mb-\alpha\le\binom m2b^2,\qquad
 0\le\alpha-b\le mb.
\]
Using \(-\ln(1-b)\le b/(1-b)\) in~\eqref{biasfull:orentropy}
yields \(\HH_o^+=O(m^2b^2)=O(\alpha^2)\).
Every child multiplier in~\eqref{biasfull:comph} is \(r\).
Consequently
\[
 \HH_\mu^+(f_m)=\HH_o^++mr\HH_\mu^+(g)
       \sim8mq^2\sim8\alpha,
\]
which proves~\eqref{biasfull:tribesratio}.
\end{proof}

\subsection{An asymmetric weighted threshold}

The preceding example is monotone but not generally a weighted threshold.
The following calculation gives a counterexample within that smaller class.
Let
\[
 h(u)=-u\ln u-(1-u)\ln(1-u)
\]
denote binary entropy in nats.

\begin{lemma}[Two-branch conditional law]\label{biasfull:mixture}
Fix \(m\ge1\), \(q,\eta\in(0,1)\), and let
\[
 a=q^m,\qquad
 q_0=\frac{(1-\eta)a}{\eta+(1-\eta)a}.
\]
Take independent \(X_0\sim\operatorname{Ber}(q_0)\) and
\(X_1,\ldots,X_m\sim\operatorname{Ber}(q)\), and set
\[
 A=\left\{mX_0+\sum_{j=1}^mX_j\ge m\right\},
 \qquad f=2\mathbf1_A-1,\qquad \alpha=\P(A).
\]
Then
\begin{align}
 \alpha&=\frac{a}{\eta+(1-\eta)a},\label{biasfull:alphamix}\\
 \TC(\mu(\cdot\mid A))
   &=m\{h(\eta+(1-\eta)q)-(1-\eta)h(q)\},
       \label{biasfull:tcmix}\\
 \frac{\I_\mu(f)}{4\alpha}
   &=1-\eta-q_0+m\eta(1-q).
       \label{biasfull:Imix}
\end{align}
\end{lemma}
\begin{proof}
The event is the disjoint union of \(\{X_0=1\}\) and
\(\{X_0=0,X_1=\cdots=X_m=1\}\).
Its probability gives~\eqref{biasfull:alphamix}.
Under the conditional law \(Q=\mu(\cdot\mid A)\),
\(\P_Q(X_0=0)=\eta\). Given \(X_0=1\), the light variables are
independent Bernoulli \(q\); given \(X_0=0\), they are all one.
Hence their marginal means are
\[
 \E_QX_0=1-\eta,\qquad
 \E_QX_j=\eta+(1-\eta)q\quad(1\le j\le m).
\]
The joint entropy is \(h(\eta)+(1-\eta)mh(q)\), while the sum of
marginal entropies is \(h(\eta)+mh(\eta+(1-\eta)q)\).
This proves~\eqref{biasfull:tcmix}.
Summing~\eqref{biasfull:pivotal} proves~\eqref{biasfull:Imix}.
\end{proof}

\begin{theorem}[Failure of weighted-threshold normalization estimates]
\label{biasfull:weightedcounter}
For every \(m\ge2\), specialize Lemma~\ref{biasfull:mixture} to
\[
 q=m^{-2},\qquad \eta=\frac1m,\qquad
 a=q^m,\qquad d=1+(m-1)a.
\]
Thus \(q_0=(m-1)a/d\) and \(\alpha=ma/d\).
For these weighted thresholds,
\begin{align}
 \frac{\alpha\TC(\mu(\cdot\mid A))}{\I_\mu(f)}
   &=\frac{\ln m+1}{8}
        +O\left(\frac{\ln m}{m}\right),\label{biasfull:weightedTC}\\
 \frac{V_\mu(f)\ln(1/V_\mu(f))
            -\sum_i\I_i^\mu(f)\ln^+(4\rho_i/V_\mu(f))}
        {\I_\mu(f)}
   &=\frac12\ln m-\ln2
        +O\left(\frac{\ln m}{m}\right),\label{biasfull:weightedcredit}\\
 \frac{\HH_\mu^+(f)}{\Delta_\mu(f)}
   &\longrightarrow2.\label{biasfull:weightedratio}
\end{align}
In particular, neither of the first two numerators is bounded above
by a universal constant times \(\I_\mu(f)\).
\end{theorem}
\begin{proof}
The light conditional mean is
\[
 r=\frac1m+\left(1-\frac1m\right)m^{-2}
     =m^{-1}+m^{-2}-m^{-3}.
\]
By~\eqref{biasfull:tcmix},
\[
 \TC(\mu(\cdot\mid A))=mh(r)-(m-1)h(q).
\]
The expansion
\(h(u)=u\ln(1/u)+u+O(u^2)\) as \(u\downarrow0\) gives
\[
 mh(r)=\ln m+1+O((\ln m)/m),\qquad
 (m-1)h(q)=O((\ln m)/m).
\]
Equation~\eqref{biasfull:Imix} is now
\begin{equation}\label{biasfull:weightedinfluence}
 \frac{\I_\mu(f)}{4\alpha}
      =2-\frac1m-q-q_0.
\end{equation}
Together these prove~\eqref{biasfull:weightedTC}.

We compute the credit residual, since the large heavy-coordinate
credit must be subtracted before the remaining cost is assessed.
The pivotal probabilities and influences satisfy
\[
 \rho_0=1-a,\qquad
 \rho_{\rm light}=(1-q_0)q^{m-1},
\]
\[
 \frac{\I_0^\mu(f)}{4\alpha}=1-\frac1m-q_0,\qquad
 \frac{\sum_{j=1}^m\I_j^\mu(f)}{4\alpha}=1-q.
\]
Since \(V:=V_\mu(f)=4\alpha(1-\alpha)\), the logarithm arguments are
\begin{equation}\label{biasfull:weightcreditargs}
 \frac{4\rho_0}{V}=\frac{1-a}{\alpha(1-\alpha)},\qquad
 \frac{4\rho_{\rm light}}{V}
      =\frac{1}{mq(1-\alpha)}=\frac{m}{1-\alpha}.
\end{equation}
Both exceed one for \(m\ge2\). Let
\[
 L=\ln(1/\alpha)=2m\ln m-\ln m+\ln d.
\]
Dividing the residual by \(4\alpha\) yields exactly
\begin{align}
 D_m:={}&\frac{V\ln(1/V)-\sum_i\I_i^\mu(f)\ln^+(4\rho_i/V)}
                    {4\alpha}\notag\\
 ={}&(1-\alpha)[L-\ln4-\ln(1-\alpha)]\notag\\
 &-\left(1-\frac1m-q_0\right)
                 [L+\ln(1-a)-\ln(1-\alpha)]\notag\\
 &-(1-q)[\ln m-\ln(1-\alpha)].
                       \label{biasfull:residualexact}
\end{align}
Here \(a=m^{-2m}\), \(q_0=O(ma)\), and \(\alpha=O(ma)\).
All terms containing \(aL\), \(q_0L\), or \(\alpha L\) tend to zero.
It follows that
\[
 D_m=\frac Lm-(1-q)\ln m-\ln4+o(1)
      =\ln m-\ln4+O((\ln m)/m).
\]
Dividing by~\eqref{biasfull:weightedinfluence} proves
\eqref{biasfull:weightedcredit}.

To prove~\eqref{biasfull:weightedratio}, let \(g\) be the sign version
of the conjunction of the \(m\) light variables. For \(S\ne\varnothing\),
\[
 P_g^\mu(S)=4a^2\left(\frac{1-q}{q}\right)^{|S|}.
\]
Summing its degree and its corrected entropy gives
\begin{align}
 \I_\mu(g)&=4ma(1-q),\notag\\
 \Delta_\mu(g)&=4a[m(1-q)-1+a],\notag\\
 \HH_\mu^+(g)&=2\Delta_\mu(g)
                 +8ma(q-a)\log(1/q).
                    \label{biasfull:andmany}
\end{align}
Consequently \(\Delta_\mu(g)\sim4ma\) and
\(\HH_\mu^+(g)\sim8ma\).
The outer OR gate has input probabilities \(q_0\) and \(a\).
Its deficit is
\[
 \Delta_o=4q_0a(1-q_0)(1-a)=O(ma^2).
\]
For an OR gate with probabilities \(x,y\), its corrected entropy is
\begin{align*}
 &8x(1-x)(1-y)^2\log\frac1{1-y}\\
 &\quad+8y(1-y)(1-x)^2\log\frac1{1-x}
       +8xy(1-x)(1-y).
\end{align*}
The inequality \(-\ln(1-u)\le u/(1-u)\) bounds this by
\((2+4/\ln2)\,4xy(1-x)(1-y)\).
Thus \(\HH_o^+=O(ma^2)\).
The heavy child is a dictator and has zero corrected entropy and
zero deficit. The light-block multiplier in
Lemma~\ref{biasfull:composition} is \(1-q_0\), so
\[
 \Delta_\mu(f)=\Delta_o+(1-q_0)\Delta_\mu(g)\sim4ma,\qquad
 \HH_\mu^+(f)=\HH_o^++(1-q_0)\HH_\mu^+(g)\sim8ma.
\]
This proves~\eqref{biasfull:weightedratio}.
\end{proof}

\begin{remark}\label{biasfull:creditremark}
Equation~\eqref{biasfull:weightedcredit} refutes precisely an estimate
with coefficient one on the pivotal credits and an additional
\(O(\I_\mu(f))\) term. It does not refute an inequality with a larger
multiplier on the credits. Replacing the fixed factor \(4\) inside
all logarithms by another fixed positive factor changes the residual
by only a bounded multiple of \(\I_\mu(f)\) for this family, and
therefore does not repair the coefficient-one estimate.
Neither~\eqref{biasfull:weightedTC} nor~\eqref{biasfull:weightedcredit}
is a counterexample to generalized FEI.
\end{remark}

\begin{proposition}[A local total-correlation obstruction]
\label{biasfull:localTC}
There is no universal \(K\) such that every increasing weighted
threshold event \(A\) has a coordinate \(j\) satisfying
\[
 \alpha\sum_{k\ne j} I_Q(X_j;X_k)\le K\I_j^\mu(2\mathbf1_A-1),
 \qquad Q=\mu(\cdot\mid A).
\]
Here the mutual informations are measured in nats.
\end{proposition}
\begin{proof}
In Lemma~\ref{biasfull:mixture}, take \(\eta=1/2\) and let
\(q\downarrow0\) with \(m\) fixed.
The conditional law converges to the equal mixture of
\((1,0,\ldots,0)\) and \((0,1,\ldots,1)\).
Hence every pairwise mutual information tends to \(\ln2\),
while~\eqref{biasfull:pivotal} gives
\(\I_j^\mu(f)/\alpha\to2\) for every \(j\).
Every proposed coordinate ratio therefore tends to \(m\ln2/2\).
Letting \(m\) increase proves the assertion.
\end{proof}

The exact recursion behind this failed proposal is
\[
 \TC(Q)=\E_{X_j}\TC(Q_{-j}\mid X_j)
                  +\sum_{k\ne j}I_Q(X_j;X_k).
\]
It follows by expanding the marginal and joint entropies.
The obstruction concerns the proposed payment of this recursion,
not the identity itself.

\subsection{What a uniform encoding of biased inputs transfers}

Let \(p_i\) be dyadic probabilities and choose arbitrary uniform-input
Boolean encoders \(g_i\) with \(\P(g_i=1)=p_i\), on disjoint blocks.
Let \(V_i=4p_i(1-p_i)\).
If ordinary uniform FEI is assumed with a constant \(C\), then
Lemma~\ref{biasfull:composition} applied to \(F=G(g_1,\ldots,g_k)\)
gives exactly
\begin{equation}\label{biasfull:encoding}
 \HH_\nu^+(G)+h_0(F)
   \le\sum_i\I_i^\nu(G)R_C(g_i),\qquad
 R_C(g)=\frac{C\I(g)-\HH^+(g)}{\Var(g)}.
\end{equation}
The relevant encoder quantity is its entropy-influence slack per
unit variance, rather than its influence alone.

\begin{proposition}[Encoder slack]\label{biasfull:encoderslack}
For the sign version \(g\) of the uniform conjunction on \(k\) bits,
with \(p=2^{-k}\),
\begin{equation}\label{biasfull:andencoder}
 R_C(g)=2+\left(\frac{C}{2(1-p)}-2\right)\log\frac1p.
\end{equation}
More generally, if uniform FEI holds with \(C_0\), then for every
\(C>C_0\) and every uniform Boolean encoder of probability
\(p\in(0,1/2]\),
\begin{equation}\label{biasfull:encodergap}
 R_C(g)\ge\frac{C-C_0}{2(1-p)}\log\frac1p.
\end{equation}
Thus no choice of encoders makes the coefficients in
\eqref{biasfull:encoding} uniformly bounded when the constant used
has a fixed positive slack over a valid FEI constant.
\end{proposition}
\begin{proof}
For the conjunction, the nonempty squared coefficients are all
\(4p^2\). Thus
\[
 \I(g)=2kp,\qquad V=4p(1-p),\qquad
 \HH^+(g)=V(2k-2),
\]
which proves~\eqref{biasfull:andencoder}.
For an arbitrary encoder,
\[
 C\I(g)-\HH^+(g)
   =C\I(g)-\HH(g)+h_0(g)\ge(C-C_0)\I(g).
\]
The uniform edge-isoperimetric inequality gives
\(\I(g)\ge2p\log(1/p)\), proving~\eqref{biasfull:encodergap}.

For completeness, this form of edge isoperimetry follows by induction.
If a coordinate's two sections have probabilities \(p_0\le p_1\),
their other influences average, and the coordinate influence is at
least \(p_1-p_0\). Writing \(p_0=p(1-t)\), \(p_1=p(1+t)\), the
induction reduces to
\[
 (1-t)\log(1-t)+(1+t)\log(1+t)\le2t,\qquad0\le t\le1.
\]
The left side is convex, with endpoint values \(0,2\), and is
therefore below the chord \(2t\).
\end{proof}

This proposition is a limitation of the inference
\eqref{biasfull:encoding}. It does not rule out a different reduction,
additional structural information about the outer gate, or the
generalized FEI conjecture itself.

\subsection{A sharp loss in a monotone dual-rail embedding}

The following calculation makes explicit the product-measure cost of the monotone encoding considered in~\cite{BrainMonotone}.

Every Boolean function on \(\{0,1\}^n\) is the restriction of a
monotone function to the antichain
\[
 x\longmapsto(x_1,1-x_1,\ldots,x_n,1-x_n).
\]
For example, take the upward closure of the encoded inputs on
which the function equals one. The following bound quantifies the
loss when the ambient law is required to be a product law.

\begin{proposition}\label{biasfull:dualrail}
Let \(Y,Z\) be independent Bernoulli bits with
\[
 \P(Y=1\mid Y\ne Z)=p,\qquad t=\P(Y\ne Z).
\]
Then
\begin{equation}\label{biasfull:railone}
 t\le\frac1{1+2\sqrt{p(1-p)}}.
\end{equation}
The bound is sharp. For independent pairs encoding independent
bits of probabilities \(p_1,\ldots,p_n\), the event \(E\) that every
pair is valid satisfies
\begin{equation}\label{biasfull:railall}
 \P(E)\le\prod_{i=1}^n
       \bigl(1+2\sqrt{p_i(1-p_i)}\bigr)^{-1}.
\end{equation}
In particular, balanced decoding forces \(\P(E)\le2^{-n}\).
\end{proposition}
\begin{proof}
The valid cell probabilities are \(pt\) and \((1-p)t\).
If the invalid cell probabilities are \(u,v\), independence implies
\[
 uv=p(1-p)t^2,\qquad u+v=1-t.
\]
Thus \((1-t)^2\ge4p(1-p)t^2\), proving~\eqref{biasfull:railone}.
Equality is attained by
\[
 \P(Y=1)=\frac{\sqrt p}{\sqrt p+\sqrt{1-p}},
 \qquad \P(Z=1)=1-\P(Y=1).
\]
Independent pairs give~\eqref{biasfull:railall} by multiplication.
\end{proof}

Conditioning on all valid pairs preserves independence of the decoded
bits, but its relative-entropy cost is
\[
 \KL(\mu(\cdot\mid E)\|\mu)=\ln\frac1{\mu(E)}
 \ge\sum_i\ln\bigl(1+2\sqrt{p_i(1-p_i)}\bigr).
\]
This is at least \(n\ln2\) for balanced decoding. The proposition
concerns this embedding; it asserts no obstruction to all conceivable
monotone embeddings.

\subsection{The gate-potential formulation on its finite-cube domain}

Let \(\mathcal Q_2\) be the dyadic probabilities in \([0,1]\).
The following formulation of the gate potential in~\cite{BrainGate}
uses only probabilities realized by finite uniform Boolean functions.
Put \(V_p=4p(1-p)\), and let \(h_2\) denote binary entropy in bits.

\begin{proposition}\label{biasfull:gatepotential}
Uniform FEI holds with constant \(C\) if and only if there is a
symmetric function \(\Gamma:\mathcal Q_2\to[0,\infty)\), with
\(\Gamma(1/2)\le C\), such that every finite Boolean gate \(G\)
and every vector of interior dyadic input probabilities satisfy
\begin{equation}\label{biasfull:gatecondition}
 \Gamma(p_{\rm out})+\HH_\nu(G)
 \le\sum_i\I_i^\nu(G)
         \frac{\Gamma(p_i)+h_2(V_{p_i})}{V_{p_i}}.
\end{equation}
Here \(p_{\rm out}=\P_\nu(G=1)\), which is dyadic. Deterministic
inputs are frozen and deleted.
\end{proposition}
\begin{proof}
For sufficiency, apply~\eqref{biasfull:gatecondition} to an arbitrary
Boolean function regarded as one gate with all input probabilities
\(1/2\). Since \(V_{1/2}=1\) and \(h_2(1)=0\),
\[
 \HH(G)+\Gamma(p_{\rm out})
      \le\Gamma(1/2)\I(G)\le C\I(G).
\]

Conversely, assume uniform FEI with \(C\) and define
\begin{equation}\label{biasfull:gammaslack}
 \Gamma_C(p)=\inf_{\substack{g\ {\rm finite\ uniform\ Boolean}\\
                              \P(g=1)=p}}
                  \{C\I(g)-\HH(g)\}.
\end{equation}
The admissible set is nonempty for dyadic \(p\).
FEI makes the infimum nonnegative; a single example makes it finite.
Output negation gives symmetry, constants give the endpoint values
zero, and a dictator gives \(\Gamma_C(1/2)\le C\).

Choose disjoint child functions of the prescribed probabilities
whose slacks approach the infima. Their raw entropy composition is
\[
 \HH(F)=\HH_\nu(G)+
    \sum_i\frac{\I_i^\nu(G)}{V_{p_i}}
                  [\HH(g_i)-h_2(V_{p_i})].
\]
This follows either from the same spectral factorization as
Lemma~\ref{biasfull:composition}, or by conditioning each child
spectrum on being nonempty. Consequently
\[
 C\I(F)-\HH(F)=
 \sum_i\frac{\I_i^\nu(G)}{V_{p_i}}
      [C\I(g_i)-\HH(g_i)+h_2(V_{p_i})]-\HH_\nu(G).
\]
The left side is at least \(\Gamma_C(p_{\rm out})\).
Letting the finitely many child slacks tend to their infima proves
\eqref{biasfull:gatecondition}.
\end{proof}

No finite uniform Boolean function has a nondyadic output probability.
Thus the infimum in~\eqref{biasfull:gammaslack} does not directly
define an all-real-bias potential; that extension would need an
additional argument. The dyadic formulation already suffices for
equivalence to FEI on all finite uniform cubes.

For comparison with the encoder bound,
\begin{equation}\label{biasfull:encoderpotential}
 \inf_{\P(g=1)=p}R_C(g)
   =\frac{\Gamma_C(p)+h_0(p)}{V_p},
 \qquad
 h_0(p)=(1-V_p)\log\frac1{1-V_p}.
\end{equation}
This identity follows from \(\HH^+(g)=\HH(g)-h_0(p)\).
If FEI holds with \(C_1<C_2\), edge isoperimetry also gives
\[
 \Gamma_{C_2}(p)\ge\Gamma_{C_1}(p)
                 +2(C_2-C_1)p\log(1/p),\qquad p\le\tfrac12.
\]

Finally, balanced gates retain the complete optimal-constant question:
\begin{equation}\label{biasfull:balancedoptimal}
 C_{\rm FEI}
  =\sup_{\substack{g\ {\rm balanced}\\\I(g)>1}}
                       \frac{\HH(g)}{\I(g)-1}.
\end{equation}
Indeed, balanced self-composition of a gate with influence \(a>1\)
and entropy \(h\) gives ratios tending to \(h/(a-1)\).
Conversely, for any nonconstant \(f\), the balanced function
\(g=z\prod_{j=1}^r f(x^{(j)})\), with a fresh uniform sign \(z\),
satisfies \(\HH(g)=r\HH(f)\) and
\(\I(g)-1=r\I(f)\). A balanced function of influence one is a
signed dictator and has zero entropy.
In particular, whenever FEI holds with \(C\),
\(\Gamma_C(1/2)=C\).

The quantifier over every gate in~\eqref{biasfull:gatecondition}
already contains unrestricted FEI at uniform input. A proof for a
restricted family of gates supplies a composition theorem for that
family; it does not remove this universal quantifier.

\section{Universal estimates that remain open}\label{sec:remaining}

The posterior forest inequality~\eqref{eq:foresttarget} is a concrete sufficient condition for unrestricted FEI. Its quantifiers matter: the representation, incidence weights, and distribution on realizable forests may depend on the function, but the final constant must not. Naming the optimal posterior does not bound its cost. The addressing calculation proves that sharing and posterior information can remove real losses; it does not prove that the remaining cost is universally bounded.

Restriction induction offers a different sufficient condition. A universal choice of a coordinate with bounded entropy increment per unit influence, or a corresponding amortized bound for a complete stopping tree, would suffice. The bent splice excludes a bound at every coordinate and the maximum-influence rule. The simultaneous reduction in Theorem~\ref{thm:simultaneous} shows that small total influence, small relative influences, near constancy, and separated Fourier supports do not collectively remove the full problem. It leaves valid existential selection and global tree accounting open.

The corrected product-measure composition identities remain useful for classes of gates. However, a claim for every Boolean outer gate of every arity at uniform input probabilities already includes the original conjecture. Likewise, a formula representation using independent occurrences cannot be transferred through the diagonal without an additional estimate. The following example makes that failure explicit.

\begin{proposition}[Independent occurrence cloning]\label{prop:cloning}
There are read-once formulas $F_m$ whose total influence and Fourier entropy both tend to zero, such that identifying occurrences of the same original variable gives $\Maj_{2m-1}$, whose Fourier entropy is at least $2$.
\end{proposition}
\begin{proof}
The positive prime DNF of $\Maj_{2m-1}$ has all
$M=\binom{2m-1}{m}$ conjunctions of $m$ positive literals. Every clause is essential: the input with exactly its $m$ coordinates positive satisfies no other clause. Replace each literal occurrence by an independent bit. The lifted formula is an OR of $M$ independent $m$-conjunctions.

Put $p=2^{-m}$ and $\mu=Mp$. Counting pivotal inputs gives
\[
 \I(F_m)=M m\,2^{1-m}(1-p)^{M-1}
        =2m\mu(1-p)^{M-1}.
\]
The central binomial coefficient is at least the average coefficient in its row, so $\mu\ge2^{m-2}/m$, whereas $\mu\le2^{m-1}$. Hence $\I(F_m)\le2m\mu e^{-\mu+p}\to0$.

The false output probability is $u_m=(1-p)^M\le e^{-\mu}$ and the dimension is $n_m=Mm$. For any Boolean function on $n$ coordinates, splitting off the empty spectral atom gives
\begin{equation}\label{eq:dimensionbound}
 \HH(f)\le h_2(\Var(f))+\Var(f)\log(2^n-1).
\end{equation}
Here $\Var(F_m)\le4e^{-\mu}$ and $n_m\le m2^{2m-1}$, so both terms tend to zero.

Identifying the copies recovers majority. This function is balanced, and for $m\ge2$ its coordinate influences are at most $1/2$. For every nonempty $S$, choose $i\in S$ and use
\[
 |\wh f(S)|=|\wh{d_i f}(S\setminus\{i\})|
 \le\E|d_i f|=\I_i(f)\le\tfrac12.
\]
Every spectral atom has probability at most $1/4$, so its spectral entropy is at least $2$.
\end{proof}

The generic minterm representation has a different limitation. Even after eliminating all latent minterm names, the uniform reference on nonempty subsets gives only~\eqref{eq:dimensionbound}, or $\HH(f)\le2\I(f)+n\Var(f)$. For balanced majority, $\Var(f)=1$ and $\I(f)\le\sqrt n$, so the last term cannot be absorbed into a universal multiple of influence. A universal forest theorem must exploit more than the existence of a finite formula.

For comparison, the general bound in~\cite[Theorem 1.3]{KKLMS20}, after taking the cutoff to be the full dimension, has the form
\[
 \HH(f)\le K\sum_S |S|\log(|S|+1)P_f(S)+K\I(f)
\]
for an absolute $K$. The class theorems discussed here remove logarithmic losses by using additional structure. Percolation supplies a spatial hierarchy whose expected occupancy costs can be summed through arm estimates~\cite{GPS10,Percolation}; ordered thresholds supply rank estimates; disjoint composition supplies a unique spectral decomposition. None of these mechanisms has yet been shown here to pay the corresponding cost for every Boolean function.

\section{Finite verification and source scope}\label{sec:verification}

The accompanying scripts verify exact identities on finite examples. They are checks of the derivations, and the proofs above do not depend on their output.

The activation and systematic-code checker visits all $256$ Boolean functions on three variables, takes two tensor copies, activates behind two controls, and applies a code with four redundant bits and distance three. It verifies $1{,}114{,}112$ Walsh coefficient identities using integer arithmetic, the exact influence and variance identities, and $3{,}840$ spectral projection injectivity checks. Entropy preservation is checked by equality of coefficient multisets. Only the separate entropy expansion check uses floating logarithms; its maximum discrepancy was less than $3\times10^{-16}$.

The auxiliary-smoothing checker treats all $16$ two-variable seeds, using rational Fourier coefficients on $26$ final coordinates. It checks the influence identities, spectral distance at least three, $5{,}616$ projection injectivity tests, and $1{,}600$ direct comparisons between the sparse Fourier expansion and the truth function. No floating arithmetic is used in these comparisons.

The multiplexer checker treats all $256$ ordered pairs of two-variable branches, including all $96$ orthogonal pairs. Entropy identities are represented as rational combinations of logarithms of primes; inequalities are reduced to integer comparisons. It also checks four addressing spectra, balanced and unbalanced subcube partitions, root conditioning, and posterior label cancellation. The restriction checker enumerates all $65{,}536$ four-variable truth tables. Walsh coefficients and influence budgets are exact dyadic quantities, while entropy rankings use floating arithmetic. Exact formulas for the reported witnesses are independently evaluated at high precision; no interval certificate for the ordering of all candidates is claimed.

\section*{Acknowledgements and funding}
The research presented in this paper was supported by the European Research
Council (ERC) under the European Union's Horizon Europe research and
innovation programme (grant agreement No.~101041711), by the Simons
Foundation as part of the Collaboration on the Mathematical and Scientific
Foundations of Deep Learning, by Heights Labs, by Convex Nexus Capital, by
the Israel Science Foundation (grant number 2258/19), and by the Israel
Science Foundation (ISF Grant 4101/25). Computational resources were
provided in part by Google Cloud through its education credits program
(credit 530550103).

\section*{Declaration of competing interest}
The author holds equity in Heights Labs and in Convex Nexus Capital and
has a management role at Convex Nexus Capital. These entities are included
in the funding declaration above.

\section*{Declaration of generative AI and AI-assisted technologies}
This paper forms part of Math Brain, an ongoing project on the use of AI
systems for research mathematics, and builds on earlier notes in that
project. In October 2026, OpenAI's GPT-6 Astra Pro, working in ChatGPT with
collaborating agents, was used to examine proof routes, derive and challenge
reductions, search the literature and the research archive, carry out finite
computations, and prepare the manuscript. Separate agents studied
restriction induction, witness forests, and biased composition, and checked
the combined constructions. These checks are not peer review. The author
directed the work and is responsible for its content and dissemination.

\section*{Data and code availability}
The source package accompanying this manuscript contains the LaTeX source,
proof notes, finite verification scripts, their recorded outputs, and a
source manifest. No experimental data were used. The package is supplied
with the manuscript.

\begin{thebibliography}{99}
\bibitem{FK96}
E.~Friedgut and G.~Kalai,
\emph{Every monotone graph property has a sharp threshold},
Proc. Amer. Math. Soc. \textbf{124} (1996), no.~10, 2993--3002.
\href{https://doi.org/10.1090/S0002-9939-96-03732-X}{doi:10.1090/S0002-9939-96-03732-X}.

\bibitem{GPS10}
C.~Garban, G.~Pete and O.~Schramm,
\emph{The Fourier spectrum of critical percolation},
Acta Math. \textbf{205} (2010), 19--104.
\href{https://doi.org/10.1007/s11511-010-0051-x}{doi:10.1007/s11511-010-0051-x}.
\href{https://arxiv.org/abs/0803.3750}{arXiv:0803.3750}.

\bibitem{GMP26}
M.~J.~Gonz\'alez, P.~MacManus and M.~C.~Pereyra,
\emph{Further evidence towards the Fourier entropy--influence conjecture},
preprint, 2026.
\href{https://arxiv.org/abs/2606.00246v2}{arXiv:2606.00246v2}.

\bibitem{KKLMS20}
E.~Kelman, G.~Kindler, N.~Lifshitz, D.~Minzer and M.~Safra,
\emph{Towards a proof of the Fourier--entropy conjecture?},
preprint, 2019, revised 2020.
\href{https://arxiv.org/abs/1911.10579v2}{arXiv:1911.10579v2}.

\bibitem{OD14}
R.~O'Donnell,
\emph{Analysis of Boolean Functions}, Cambridge University Press, 2014.
\href{https://www.cs.cmu.edu/~odonnell/papers/Analysis-of-Boolean-Functions-by-Ryan-ODonnell.pdf}{Author's edition}.

\bibitem{ODT13}
R.~O'Donnell and L.-Y.~Tan,
\emph{A composition theorem for the Fourier entropy--influence conjecture},
in: ICALP 2013, Lecture Notes in Comput. Sci. \textbf{7965}, Springer, 2013, 780--791.
\href{https://doi.org/10.1007/978-3-642-39206-1_66}{doi:10.1007/978-3-642-39206-1\_66}.
\href{https://arxiv.org/abs/1304.1347}{arXiv:1304.1347}.

\bibitem{OWZ11}
R.~O'Donnell, J.~Wright and Y.~Zhou,
\emph{The Fourier entropy--influence conjecture for certain classes of Boolean functions},
in: ICALP 2011, Lecture Notes in Comput. Sci. \textbf{6755}, Springer, 2011, 330--341.
\href{https://doi.org/10.1007/978-3-642-22006-7_28}{doi:10.1007/978-3-642-22006-7\_28}.

\bibitem{Percolation}
Y.~Shmalo,
\emph{Entropy and singularity of spectral and pivotal sets in planar percolation},
preprint, 2026.
\href{https://arxiv.org/abs/2609.15388}{arXiv:2609.15388}.
Only its entropy and spatial encoding results are used for comparison here.

\bibitem{Threshold}
Y.~Shmalo,
\emph{Fourier entropy of threshold functions and their compositions},
companion manuscript, 11 October 2026.

\bibitem{Shalev18}
G.~Shalev,
\emph{On the Fourier entropy influence conjecture for extremal classes},
preprint, 2018, revised 2019.
\href{https://arxiv.org/abs/1806.03646v2}{arXiv:1806.03646v2}.

\bibitem{BrainErasure}
Math Brain research notes,
\emph{FEI: a corrected reduction and constant-preserving erasure obstructions},
18 September 2026, unpublished.
Repository path: \texttt{research\_pushes/2026-09-18-sidorenko-fei-ising/FEI.md}.

\bibitem{BrainCritical}
Math Brain research notes,
\emph{FEI continuation: global attempts and bounded critical-core search},
30 September 2026, unpublished.
Repository path: \texttt{research\_pushes/2026-09-30-fei-critical-search/ATTEMPTS.md}.

\bibitem{BrainAllocation}
Math Brain research notes,
\emph{FEI logarithmic allocation report}, 30 September 2026, unpublished.
Repository path: \texttt{research\_pushes/2026-09-30-fei-log-allocation/REPORT.md}.

\bibitem{BrainGate}
Math Brain research notes,
\emph{FEI gate reformulation}, 22 September 2026, unpublished.
Repository continuation: \texttt{proofnet24/P03-fei/continuations/2026-09-22-gate-reformulation}.

\bibitem{BrainMonotone}
Math Brain research notes,
\emph{Monotone and full-monotone FEI continuations},
24 September 2026, unpublished.
Repository continuations: \texttt{2026-09-24-monotone} and \texttt{2026-09-24-full-monotone}.
\end{thebibliography}

\end{document}
